\begin{filecontents*}{synthesis.bib}
@article{mashru2023,
  title = {A Decentralised Instant Messaging Application with End-to-End Encryption},
  author = {Dev Mashru and Ganesha Maruthi Mangipudi and Harivardhan Swamy and Shivakumar Halangali and Sushma E},
  year = {2023},
  journal = {2023 20th Learning and Technology Conference (L\&T)},
  doi = {10.1109/LT58159.2023.10092319},
  url = {https://ieeexplore.ieee.org/document/10092319/},
}

@article{alsop2025,
  title = {(Poster) ARSecure: A Novel End-to-End Encryption Messaging System using Augmented Reality},
  author = {Hamish Alsop and Douglas Alsop and Joseph Solomon and Liam Aumento and Mark Butters and Cameron Millar and Leandros Maglaras and Yagmur Yigit and Berk Canberk and Naghmeh Moradpoor},
  year = {2025},
  journal = {2025 21st International Conference on Distributed Computing in Smart Systems and the Internet of Things (DCOSS-IoT)},
  doi = {10.1109/DCOSS-IoT65416.2025.00046},
  url = {https://ieeexplore.ieee.org/document/11096238/},
}

@article{p2023,
  title = {Hybrid Encryption through End to End in Messaging Service Applications},
  author = {Vidya Sagar P and Dhinesh K S and Jayakumar M and Hemamalini D},
  year = {2023},
  journal = {2023 International Conference on Inventive Computation Technologies (ICICT)},
  doi = {10.1109/ICICT57646.2023.10134120},
  url = {https://ieeexplore.ieee.org/document/10134120/},
}

@article{jain2021,
  title = {Vulnerability Analysis of a Signal-based Messenger},
  author = {Kartikeya Jain and Anagha Ananth and Prasad Honnavalli},
  year = {2021},
  journal = {2021 IEEE Bombay Section Signature Conference (IBSSC)},
  doi = {10.1109/IBSSC53889.2021.9673482},
  url = {https://ieeexplore.ieee.org/document/9673482/},
}

@article{kamkuemah2021,
  title = {Reasoning about Authentication and Secrecy in the Signal Protocol},
  author = {Martha N. Kamkuemah},
  year = {2021},
  journal = {2021 International Conference on Electrical, Computer and Energy Technologies (ICECET)},
  doi = {10.1109/ICECET52533.2021.9698415},
  url = {https://ieeexplore.ieee.org/document/9698415/},
}

@article{morales2024,
  title = {Towards the Removal of Identification and Authentication Authority from IM Systems},
  author = {David A. Cordova Morales and Ahmad Samer Wazan and Romain Laborde and Muhammad Imran Taj and Adib Habbal and Gina Gallegos-García},
  year = {2024},
  journal = {2024 8th Cyber Security in Networking Conference (CSNet)},
  doi = {10.1109/CSNet64211.2024.10851747},
  url = {https://ieeexplore.ieee.org/document/10851747/},
}

@article{boras2025,
  title = {Post-Quantum Cryptography in Secure Instant Messaging Protocols},
  author = {Katarina Boras},
  year = {2025},
  journal = {2025 MIPRO 48th ICT and Electronics Convention},
  doi = {10.1109/MIPRO65660.2025.11132015},
  url = {https://ieeexplore.ieee.org/document/11132015/},
}

@article{krishnapriya2021,
  title = {Forensic Extraction and Analysis of Signal Application in Android Phones},
  author = {S Krishnapriya and V S Priyanka and S Satheesh Kumar},
  year = {2021},
  journal = {2021 International Conference on Forensics, Analytics, Big Data, Security (FABS)},
  doi = {10.1109/FABS52071.2021.9702702},
  url = {https://ieeexplore.ieee.org/document/9702702/},
}

@article{bobrysheva2021,
  title = {Post-quantum Secure Group Messaging},
  author = {Julia Bobrysheva and Sergey Zapechnikov},
  year = {2021},
  journal = {2021 IEEE Conference of Russian Young Researchers in Electrical and Electronic Engineering (ElConRus)},
  doi = {10.1109/ElConRus51938.2021.9396513},
  url = {https://ieeexplore.ieee.org/document/9396513/},
}

@article{erdenebaatar2023,
  title = {Instant Messaging Application Encrypted Traffic Generation System},
  author = {Zolboo Erdenebaatar and Biswajit Nandy and Nabil Seddigh and Riyad Alshammari and Marwa Elsayed and Nur Zincir-Heywood},
  year = {2023},
  journal = {NOMS 2023-2023 IEEE/IFIP Network Operations and Management Symposium},
  doi = {10.1109/NOMS56928.2023.10154297},
  url = {https://ieeexplore.ieee.org/document/10154297/},
}

@article{sarvaiya2025,
  title = {Detecting Presence Of Malicious Hub in MIMI Protocol for Cross-Platform Messaging Interoperability},
  author = {Harditya Sarvaiya and Eric W. Burger},
  year = {2025},
  journal = {MILCOM 2025 - 2025 IEEE Military Communications Conference (MILCOM)},
  doi = {10.1109/MILCOM64451.2025.11310273},
  url = {https://ieeexplore.ieee.org/document/11310273/},
}

@article{nelson2024,
  title = {Metadata Privacy Beyond Tunneling for Instant Messaging},
  author = {Boel Nelson and Elena Pagnin and Aslan Askarov},
  year = {2024},
  journal = {2024 IEEE 9th European Symposium on Security and Privacy (EuroS\&P)},
  doi = {10.1109/EuroSP60621.2024.00044},
  url = {https://ieeexplore.ieee.org/document/10629009/},
}

@article{m2024,
  title = {End to End Encrypted Messaging Application Using ORDEX},
  author = {Vishwas M and V Raghav Loknath and Vinti Agrawal and Vishwa Mehul Mehta and Sarasvathi V},
  year = {2024},
  journal = {2024 6th International Conference on Electrical, Control and Instrumentation Engineering (ICECIE)},
  doi = {10.1109/ICECIE63774.2024.10815684},
  url = {https://ieeexplore.ieee.org/document/10815684/},
}

@article{bozorgi2023,
  title = {I Still Know What You Did Last Summer: Inferring Sensitive User Activities on Messaging Applications Through Traffic Analysis},
  author = {Ardavan Bozorgi and Alireza Bahramali and Fateme Rezaei and Amirhossein Ghafari and Amir Houmansadr and Ramin Soltani and Dennis Goeckel and Don Towsley},
  year = {2023},
  journal = {IEEE Transactions on Dependable and Secure Computing},
  doi = {10.1109/TDSC.2022.3218191},
  url = {https://ieeexplore.ieee.org/document/9933634/},
}

@article{basem2023,
  title = {Stick: An End-to-End Encryption Protocol Tailored for Social Network Platforms},
  author = {Omar Basem and Abrar Ullah and Hani Ragab Hassen},
  year = {2023},
  journal = {IEEE Transactions on Dependable and Secure Computing},
  doi = {10.1109/TDSC.2022.3152256},
  url = {https://ieeexplore.ieee.org/document/9716793/},
}

@article{r2025,
  title = {EnConvo: Secure End-to-End Encrypted Messaging Application},
  author = {Muthu Pandeeswari R and Priya Dharshini and Sajay Prakash K},
  year = {2025},
  journal = {2025 International Conference on Electronics and Renewable Systems (ICEARS)},
  doi = {10.1109/ICEARS64219.2025.10940216},
  url = {https://ieeexplore.ieee.org/document/10940216/},
}

@article{kanimozhi2025,
  title = {Decentralized Social Networking with Artificial Intelligence for Enhanced Connection and Updates},
  author = {S. Kanimozhi and Anuj S and Santheep W S and Hari Prasath S},
  year = {2025},
  journal = {2025 5th Asian Conference on Innovation in Technology (ASIANCON)},
  doi = {10.1109/ASIANCON66527.2025.11281145},
  url = {https://ieeexplore.ieee.org/document/11281145/},
}

@article{ayed2025,
  title = {Developed a Multiple Ratchet Group Key Agreement Protocol},
  author = {Khalid J. Ayed and Omar A. Dawood},
  year = {2025},
  journal = {2025 4th International Conference on Computing and Information Technology (ICCIT)},
  doi = {10.1109/ICCIT63348.2025.10989424},
  url = {https://ieeexplore.ieee.org/document/10989424/},
}

@article{thapa2025,
  title = {Efficiency in Chat Application Encryption: A Comparative Review with Proposed Enhancements},
  author = {Jikesh Thapa and MD Nashid Anjum and Rashid Khokhar},
  year = {2025},
  journal = {2025 IEEE/ACIS 29th International Conference on Software Engineering, Artificial Intelligence, Networking and Parallel/Distributed Computing (SNPD)},
  doi = {10.1109/SNPD65828.2025.11253003},
  url = {https://ieeexplore.ieee.org/document/11253003/},
}

@article{angom2025,
  title = {MLXDH: Security Hardening of Signal’s Initial Key Establishment Using ML-KEM and ML-DSA},
  author = {Akash Angom and Nirmalya Kar and Tribid Debbarma and Priyanka Biswas},
  year = {2025},
  journal = {2025 IEEE 6th India Council International Subsections Conference (INDISCON)},
  doi = {10.1109/INDISCON66021.2025.11254562},
  url = {https://ieeexplore.ieee.org/document/11254562/},
}

@article{zhao2022,
  title = {5G Messaging: System Insecurity and Defenses},
  author = {Jinghao Zhao and Qianru Li and Zengwen Yuan and Zhehui Zhang and Songwu Lu},
  year = {2022},
  journal = {2022 IEEE Conference on Communications and Network Security (CNS)},
  doi = {10.1109/CNS56114.2022.9947238},
  url = {https://ieeexplore.ieee.org/document/9947238/},
}

@article{sharma2023,
  title = {A Post-Quantum End-to-End Encryption Protocol},
  author = {Shashvat Sharma and Meenakshi Tripathi and Harish Kumar Sahu and Ashish Karan},
  year = {2023},
  journal = {2023 IEEE International Conference on Advanced Networks and Telecommunications Systems (ANTS)},
  doi = {10.1109/ANTS59832.2023.10469296},
  url = {https://ieeexplore.ieee.org/document/10469296/},
}

@article{yi2020,
  title = {eDirect: Energy-Efficient D2D-Assisted Relaying Framework for Cellular Signaling Reduction},
  author = {Xiaomeng Yi and Li Pan and Yanqi Jin and Fangming Liu and Minghua Chen},
  year = {2020},
  journal = {IEEE/ACM Transactions on Networking},
  doi = {10.1109/TNET.2020.2973540},
  url = {https://ieeexplore.ieee.org/document/9018118/},
}

@article{oppliger2024,
  title = {Signal and Messaging Layer Security},
  author = {Rolf Oppliger},
  year = {2024},
  journal = {Signal and Messaging Layer Security},
  url = {https://ieeexplore.ieee.org/document/11076157/},
}

@article{cao2022,
  title = {Analysis on end-to-end transmission protocol and its work performance},
  author = {Yuchen Cao and Tianyi Zhang},
  year = {2022},
  journal = {2022 International Conference on Data Analytics, Computing and Artificial Intelligence (ICDACAI)},
  doi = {10.1109/ICDACAI57211.2022.00049},
  url = {https://ieeexplore.ieee.org/document/9973835/},
}

@article{jiang2024,
  title = {Purified Authorization Service With Encrypted Message Moderation},
  author = {Peng Jiang and Qi Liu and Liehuang Zhu},
  year = {2024},
  journal = {IEEE Transactions on Information Forensics and Security},
  doi = {10.1109/TIFS.2024.3393391},
  url = {https://ieeexplore.ieee.org/document/10508244/},
}

@article{vladimirovich2020,
  title = {Formalization of Organization of Signaling Protocol While Making Online Calls Using Automata Theory},
  author = {Pashchenko Dmitry Vladimirovich and Mitrokhin Maxim Aleksandrovich and Trokoz Dmitry Anatolyevich and Sinev Mikhail Petrovich and Savateev Maxim Valerievich and Iskhakov Nikita Valerievich and Rodionov Vladislav Sergeevich},
  year = {2020},
  journal = {2020 Moscow Workshop on Electronic and Networking Technologies (MWENT)},
  doi = {10.1109/MWENT47943.2020.9067350},
  url = {https://ieeexplore.ieee.org/document/9067350/},
}

@article{brigham2025,
  title = {No Safety in Numbers: Traffic Analysis of Sealed-Sender Groups in Signal},
  author = {Eric Brigham and Nicholas Hopper},
  year = {2025},
  journal = {2025 22nd Annual International Conference on Privacy, Security, and Trust (PST)},
  doi = {10.1109/PST65910.2025.11268871},
  url = {https://ieeexplore.ieee.org/document/11268871/},
}

@article{dowling2021,
  title = {Secure Messaging Authentication against Active Man-in-the-Middle Attacks},
  author = {Benjamin Dowling and Britta Hale},
  year = {2021},
  journal = {2021 IEEE European Symposium on Security and Privacy (EuroS\&P)},
  doi = {10.1109/EuroSP51992.2021.00015},
  url = {https://ieeexplore.ieee.org/document/9581268/},
}

@article{oppliger2020,
  title = {End-to-End Encrypted Messaging},
  author = {Rolf Oppliger},
  year = {2020},
  journal = {End-to-End Encrypted Messaging},
  url = {https://ieeexplore.ieee.org/document/9118792/},
}

@article{samper2025,
  title = {SoK: Self-Generated Nudes over Private Chats: How can Technology Contribute to a Safer Sexting?},
  author = {Joel Samper and Bernardo Ferreira},
  year = {2025},
  journal = {2025 IEEE Symposium on Security and Privacy (SP)},
  doi = {10.1109/SP61157.2025.00234},
  url = {https://ieeexplore.ieee.org/document/11023422/},
}
\end{filecontents*}

\documentclass{article}
\usepackage[utf8]{inputenc}
\usepackage{hyperref}
\usepackage{enumitem}
\usepackage{longtable}
\usepackage{booktabs}
\usepackage{array}
\usepackage[backend=biber,style=authoryear]{biblatex}
\addbibresource{synthesis.bib}

\title{Literature Synthesis}
\date{2026-01-15}

\begin{document}
\maketitle

\section{Research Question}

\subsection{Search Parameters}

\begin{description}
\item[Query] What research has been done releated to instant messaging and the Signal protocol since 2020?
\item[Providers] ieee
\end{description}

\subsection{Research Context}
What research has been done releated to instant messaging and the Signal
protocol since 2020?

\subsection{Summary}

\begin{description}
\item[Date] 2026-01-15 11:56
\item[Total Papers] 100
\item[Kept] 31
\item[Discarded] 69
\item[Pending] 0
\end{description}


\section{Synthesis}

\subsection{android}
Android-focused work since 2020 has examined Signal from the standpoint of post-hoc evidence recovery rather than protocol redesign, positioning the platform as a locus where end-to-end encryption coexists with residual artefacts useful for investigations \cite{krishnapriya2021}. Krishnapriya et al. explicitly frame Signal’s privacy-preserving reputation as driving adoption away from other mainstream messengers, which in turn increases the need for forensic methods tailored to Signal on Android \cite{krishnapriya2021}. Within this framing, the central research contribution is methodological: they propose acquiring Signal application data by decrypting the Signal backup file and then recovering chat messages and attachments from Android devices \cite{krishnapriya2021}. This emphasis highlights an Android-specific research direction in which security properties at the transport/protocol layer are treated as given, while practical access to user data is pursued via device-side storage and backup mechanisms \cite{krishnapriya2021}.

\subsection{authentication}
Research since 2020 has increasingly treated authentication in Signal-style instant messaging not as a settled byproduct of end-to-end encryption, but as a distinct problem spanning formal properties, human-mediated key verification, and system-level trust assumptions. Two complementary formal-methods threads interrogate what it means for users to be authenticated in a ratcheted messaging setting and whether Signal’s deployed mechanisms meet those notions \cite{dowling2021,kamkuemah2021}, while a third line argues that persistent authentication weaknesses are rooted in centralized identification and public-key infrastructure and proposes moving key authentication outside the IM provider altogether \cite{morales2024}.

A key point of divergence is the abstraction layer at which authentication is modeled. Kamkuemah frames authentication and secrecy in epistemic terms—reasoning explicitly about ``who knows what''—to express intended security properties prior to implementation and then analyze Signal’s behavior against those properties \cite{kamkuemah2021}. In contrast, Dowling and Hale treat authentication failures as arising from the interaction between cryptographic key exchange and user-mediated out-of-band (OOB) verification, proposing a formal model that captures how real users attest to long-term public keys and how this composes with a ratcheted key exchange \cite{dowling2021}. These different formalisms yield different diagnostic power: epistemic-logic analysis supports claims that Signal achieves authentication in the modeled knowledge-based sense \cite{kamkuemah2021}, whereas the user-mediated authentication model highlights an active, post-compromise man-in-the-middle setting in which the Signal application’s authentication procedure cannot be proven secure \cite{dowling2021}.

Across these works, the threat model around active attackers and compromise is central to how authentication is evaluated. Dowling and Hale emphasize that although ratcheted protocols can offer confidentiality recovery properties after compromise (e.g., PCS), authenticity is generally lost against active attackers, motivating a definition of active post-compromise entity authentication security and an analysis showing Signal’s user-mediated authentication falls short under that definition \cite{dowling2021}. Kamkuemah similarly foregrounds adversarial knowledge of past and future keys following exposure of one key as part of Signal’s promised protections, but evaluates authentication via epistemic security properties that conclude the protocol remains secure \cite{kamkuemah2021}. Read together, these papers suggest that ``Signal is secure'' can be simultaneously true in a protocol-knowledge sense and insufficient in a deployment-relevant sense when user-mediated verification and post-compromise active attacks are treated as first-class concerns \cite{dowling2021,kamkuemah2021}.

A second synthesis dimension is where authentication authority is located and how much users must shoulder. Dowling and Hale analyze protocols that reduce reliance on trusted third parties by shifting public-key attestation to users via OOB verification, then show that this very reliance on user mediation leaves gaps against active MitM after compromise unless the mechanism is strengthened, offering a straightforward fix to enable detection of an active adversary in Signal’s setting \cite{dowling2021}. Morales et al.\ approach the same trust bottleneck from the opposite direction: rather than relying on IM providers (or implicitly on users alone) for public-key authentication, they propose a decentralized framework based on Trust over IP using Decentralized Identifiers and Verifiable Credentials to manage public keys outside provider control, explicitly aiming to eliminate IM systems as identification and authentication authorities \cite{morales2024}. This decentralization proposal aligns with the diagnosis that public-key authentication remains problematic despite the availability of Signal-like encryption, attributing weaknesses to centralized infrastructure, trust issues, and user misunderstanding of E2EE \cite{morales2024}, and implicitly complements formal results by relocating the authentication substrate that user-mediated models identify as fragile \cite{dowling2021}.

Overall, post-2020 work on authentication around Signal and secure messaging coalesces around three interconnected insights: (i) authentication requires explicit formalization beyond confidentiality, especially under active and post-compromise conditions \cite{dowling2021}; (ii) the choice of logic/model (knowledge-centric vs.\ user-mediated compositional modeling) materially affects whether Signal’s authentication is judged sufficient \cite{kamkuemah2021,dowling2021}; and (iii) architectural changes to key authentication (e.g., decentralized identifiers/credentials) are proposed as a system-level remedy to persistent trust and usability weaknesses in public-key authentication for IM, even when Signal-like encryption is assumed \cite{morales2024}.

\subsection{client-side-scanning}

Work since 2020 that engages with client-side scanning (CSS) in the context of end-to-end encrypted instant messaging emphasizes CSS as a direct erosion of E2EE guarantees by moving content inspection onto endpoints rather than intermediaries \cite{alsop2025}. In this framing, the central technical tension is not primarily the cryptography in transit, but the trustworthiness of the client device that performs encryption/decryption and can be compelled or modified to scan message payloads (text and media) on send/receive \cite{alsop2025}. 

Alsop et al.\ propose a design response that treats CSS as an endpoint compromise problem and attempts to reduce the attack surface by relocating encryption and decryption away from the general-purpose messaging client \cite{alsop2025}. Their ARSecure system uses augmented reality glasses as an externalized, security-oriented component to handle cryptographic operations, with the explicit goal of mitigating client-side vulnerabilities that enable CSS and thereby restoring a stronger interpretation of “tamper-proof, private communication” for sensitive environments \cite{alsop2025}. Conceptually, this approach contrasts with defenses that would aim to detect or regulate scanning within the client software itself; instead, ARSecure seeks to make scanning harder by separating where plaintext is rendered and where cryptographic transformations occur, so the conventional smartphone/desktop client is less able to access or instrument message content in a way that supports CSS \cite{alsop2025}. 

Across the CSS theme, this positions hardware-assisted or out-of-band interfaces as a privacy-preserving countermeasure to policy- or platform-driven scanning demands: rather than negotiating what gets scanned, ARSecure attempts to structurally limit the client’s capability to scan at all by shifting trust to a dedicated external device \cite{alsop2025}. At the same time, the work implicitly underscores that debates about CSS cannot be resolved solely by selecting a strong protocol (e.g., E2EE) if the endpoint remains a universal execution environment that can be repurposed for surveillance functions \cite{alsop2025}.

\subsection{data-extraction}
Recent work on Signal-related instant messaging forensics has concentrated on extracting user content from Android devices despite end-to-end encryption, by targeting local artefacts rather than attempting to break transport security \cite{krishnapriya2021}. In this line, Krishnapriya et al.\ frame data-extraction as a problem of acquiring and interpreting on-device application data, proposing a methodology centred on decrypting the Signal backup file to recover chat messages and associated attachments \cite{krishnapriya2021}. This approach implicitly treats the backup as a privileged evidentiary source: while Signal’s communications are designed to be secure in transit, the paper’s contribution is to shift forensic access to a user-enabled persistence mechanism and then operationalize its decryption for evidentiary recovery \cite{krishnapriya2021}. Within the theme of data-extraction, the work thus exemplifies a practical, acquisition-first pathway—identifying an accessible at-rest data container (the backup), performing decryption, and extracting higher-level conversational artefacts (messages and attachments) suitable for analysis \cite{krishnapriya2021}.

\subsection{dataset-generation}

Recent work on dataset-generation for encrypted instant messaging traffic emphasizes automation and ground-truth labeling as prerequisites for downstream analysis, motivated by the scarcity of publicly available datasets for IMAs \cite{erdenebaatar2023}. Erdenebaatar et al. design an encrypted traffic generation system that operationalizes this need by automatically producing and labeling traffic for multiple IMAs—including Signal alongside Discord, Facebook Messenger, Microsoft Teams, Telegram, and WhatsApp—thereby positioning dataset-generation as a cross-application infrastructure problem rather than an app-specific data collection task \cite{erdenebaatar2023}.

A key methodological choice in this system is to generate data by emulating user behavior while capturing, filtering, and labeling traffic directly on an Android device \cite{erdenebaatar2023}. This end-to-end, on-device pipeline links activity generation with capture and labeling, suggesting a practical route to reducing label noise and improving reproducibility compared with ad hoc, manually driven data collection workflows \cite{erdenebaatar2023}. By framing the contribution as a companion to prior flow-based encrypted IMA traffic analysis, the paper also connects dataset-generation tightly to analytics requirements: the generated labeled flows are intended to support data-driven methods in encrypted traffic analysis, including for Signal traffic \cite{erdenebaatar2023}.

\subsection{decentralized-architecture}
Recent post-2020 work on decentralized architecture for instant messaging emphasizes reducing reliance on corporate-controlled centralized infrastructure while retaining end-to-end confidentiality. Across the two papers, decentralization is framed both as a resilience measure against single points of failure and as a governance shift toward user-controlled data ownership \cite{mashru2023,kanimozhi2025}. While both approaches retain end-to-end encryption as a baseline privacy mechanism, they diverge in where state resides and how availability is maintained under decentralization \cite{mashru2023,kanimozhi2025}.

A key architectural contrast concerns data placement and persistence. Mashru et al.\ propose a decentralized messaging application “hosted on the cloud,” motivated by avoiding centralized-server dependence and improving fault tolerance \cite{mashru2023}. In comparison, Kanimozhi et al.\ center decentralization on local-first storage—profiles and content remain on users’ devices—while introducing ephemeral caching servers that temporarily mirror data on decentralized nodes during active interactions and erase it after the session \cite{kanimozhi2025}. This juxtaposition highlights two distinct decentralization strategies: decentralizing service operation for fault tolerance via cloud-hosted components \cite{mashru2023} versus decentralizing data custody by making the user device the primary store and treating network infrastructure as transient support \cite{kanimozhi2025}.

Both papers position decentralization as intertwined with broader platform goals beyond message confidentiality. Mashru et al.\ primarily treat decentralization as a remedy for market domination and the central point-of-failure problem, with end-to-end encryption providing privacy assurances within that architecture \cite{mashru2023}. Kanimozhi et al.\ similarly oppose exploitative centralized platforms, but embed decentralization into a wider social-networking stack that blends mainstream app interaction patterns with end-to-end encryption “similar to Signal,” alongside AI-driven personalization and a network-wide trust/reputation mechanism \cite{kanimozhi2025}. In this sense, decentralization is not only an infrastructural choice but also an enabling substrate for additional governance and interaction features (e.g., reputation-based misuse deterrence) that operate across distributed nodes \cite{kanimozhi2025}.

Taken together, the two papers suggest that since 2020, decentralized-architecture research in instant messaging and related social communication systems has explored (i) decentralizing operational infrastructure to reduce single points of failure while maintaining end-to-end encrypted messaging \cite{mashru2023}, and (ii) decentralizing data ownership via local storage complemented by ephemeral, session-scoped caching in the network, paired with Signal-like end-to-end encryption and additional AI-augmented social features \cite{kanimozhi2025}. The common throughline is the use of decentralization to rebalance control and resilience, with architectural differences centered on whether the system’s “center of gravity” is distributed cloud operation or device-centric data custody with transient network assistance \cite{mashru2023,kanimozhi2025}.

\subsection{decentralized-identities}
Research since 2020 linking instant messaging security to the Signal protocol has increasingly treated cryptographic confidentiality as comparatively mature while reframing user identity and public-key authentication as the remaining structural vulnerability, particularly where providers implicitly act as identification and authentication authorities \cite{morales2024}. Within this framing, Morales et al.\ propose decoupling authentication from the IM service by moving public-key authentication into a decentralized identity layer, explicitly targeting the centralized trust assumptions and user-facing misunderstandings that persist even when end-to-end encryption is provided via Signal-like designs \cite{morales2024}.

A key contribution is the articulation of a provider-agnostic authentication architecture grounded in the Trust over IP (ToIP) model, where Decentralized Identifiers (DIDs) and Verifiable Credentials (VCs) are used to bind and manage users' public keys outside the IM provider’s infrastructure \cite{morales2024}. This shifts the locus of trust from platform-specific public key infrastructure to decentralized, transparent mechanisms, with the intended effect of eliminating the provider’s role as an ``artificial'' authority over identification and authentication while preserving secure messaging goals \cite{morales2024}. In synthesis, the work positions decentralized identities not as an alternative to Signal’s encryption, but as an orthogonal layer designed to address the residual authentication problem by redistributing control over key assertions and identity proofs away from centralized IM operators \cite{morales2024}. Finally, the proposal ties this architectural shift to interoperability and compliance motivations, arguing alignment with regulatory frameworks such as the Digital Markets Act (DMA) and eIDAS as an additional driver for decentralized identity-based key authentication in IM ecosystems \cite{morales2024}.

\subsection{double-ratchet}
Since 2020, work engaging the Signal family of mechanisms has used the double-ratchet primarily as a security baseline to be extended or adapted, rather than re-litigated as a one-to-one primitive, with the most visible variations arising when targeting either post-quantum assumptions or group messaging dynamics \cite{bobrysheva2021,ayed2025,m2024}. Across the provided papers, the double-ratchet is treated as the component responsible for ongoing message-key evolution (and associated forward-secrecy style properties), while surrounding key agreement and group-management layers are redesigned to meet new constraints such as quantum resistance, scalability, or performance \cite{bobrysheva2021,ayed2025,m2024}.

A key line of divergence concerns how the double-ratchet is “lifted” beyond pairwise sessions. Bobrysheva and Zapechnikov explicitly frame their proposal as an \emph{extended Double Ratchet} for group contexts and tie that extension to a new group key establishment scheme based on isogenies of elliptic curves, motivated by post-quantum resistance \cite{bobrysheva2021}. In contrast, Ayed and Dawood aim to preserve “the security features provided by the double ratchet protocol in one-on-one communication” while engineering a scalable group key agreement, proposing a multiple ratchet group key agreement (MRGKA) that operationalizes group dynamics via partitioning into bounded-size subgroups (up to 15 members) that behave “as if each subgroup were a single user entity” for intra-group privacy \cite{ayed2025}. Read together, these approaches reflect two distinct strategies for group adaptation: (i) redefining the underlying group key establishment primitive and then describing it as an extension of double-ratchet behavior \cite{bobrysheva2021}, versus (ii) structuring the group to retain double-ratchet-like properties by decomposing it into manageable entities and adding protocol machinery for membership churn \cite{ayed2025}.

The papers also differ on what they optimize around the double-ratchet. M et al.\ embed Double Ratchet alongside X3DH in an application design but focus their novelty on accelerating the key-exchange step via ORDEX, claiming substantially higher speed than ECDH-based exchange while keeping the rest of the pipeline consistent with a layered E2EE design (AES for message encryption, plus X3DH and Double Ratchet) \cite{m2024}. This positions double-ratchet as a stable “known-good” component whose surrounding handshake can be swapped for performance, in contrast to the group-oriented papers where the principal challenge is maintaining ratcheting-derived security properties under asynchronous, dynamic membership changes \cite{ayed2025,bobrysheva2021}. Ayed and Dawood further specify a concrete cryptographic suite (Chacha20/Poly1305, SHA256, XEdDSA, Curve2559) and introduce an Invite Authentication Code (IAC) to strengthen subsequent sessions in their symmetric ratchet approach, indicating that practical group ratcheting demands additional authentication and lifecycle controls beyond the core ratchet update mechanism \cite{ayed2025}.

Finally, the collection highlights that “double-ratchet” research since 2020 is not monolithic: it is simultaneously a label for (a) post-quantum re-grounding of the key establishment feeding a ratchet in group messaging \cite{bobrysheva2021}, (b) architectural decomposition (subgrouping) and auxiliary controls (IAC) to make ratchet-like guarantees workable at group scale \cite{ayed2025}, and (c) application-driven integration where the ratchet remains constant while the initial exchange is optimized for speed \cite{m2024}. The commonality is the preservation of a ratcheting state-evolution concept for ongoing confidentiality, while the surrounding protocol layers are where innovation concentrates under new threat models and deployment constraints \cite{bobrysheva2021,ayed2025,m2024}.

\subsection{end-to-end-encryption}

Since 2020, end-to-end encryption (E2EE) research in instant messaging has simultaneously consolidated around the Signal protocol as a de facto baseline and diversified into new architectural and governance directions. Oppliger positions the Signal protocol as central to contemporary E2EE messaging and highlights that privacy risks persist beyond payload confidentiality, including traffic-analysis-relevant metadata and features such as self-destructing messages \cite{oppliger2020}. Against this baseline, later work spans (i) deploying E2EE in alternative system architectures, (ii) strengthening E2EE against emerging threats (notably quantum adversaries and client-side compromise), (iii) addressing the moderation/accountability tension created by encrypted content, and (iv) evaluating the efficiency costs of E2EE on commodity devices.

A first line of work treats E2EE as necessary but insufficient without changing the surrounding infrastructure. Mashru et al.\ argue that reliance on centralized servers creates single points of failure and corporate control, motivating a decentralized instant messaging application that still provides E2EE for message privacy \cite{mashru2023}. Kanimozhi et al.\ similarly combine decentralized social networking with E2EE “similar to Signal,” aligning privacy goals with user-controlled data storage and ephemeral caching to improve accessibility during interactions \cite{kanimozhi2025}. Read together, these systems-oriented contributions frame E2EE as a portable security primitive that can be embedded into architectures aimed at resilience and user autonomy, rather than being exclusive to large, centrally operated messengers \cite{mashru2023,kanimozhi2025}. In this framing, E2EE’s role is stable (confidentiality between endpoints), while the surrounding system design varies to reduce dependence on central intermediaries \cite{mashru2023,kanimozhi2025}.

A second line of work reacts to the prospect that classical public-key components underpinning popular E2EE could become obsolete under quantum attack. Sharma et al.\ explicitly identify the vulnerability of ECDH-based key exchange (as used in the Signal protocol) to quantum capabilities and motivate replacing it with quantum-safe key exchange while preserving the broader security properties associated with Signal-based secure messaging \cite{sharma2023}. In parallel, EnConvo proposes a post-quantum-capable E2EE application design that combines Kyber-1024 with X25519 for compatibility, adds hardware-backed key protection (e.g., trusted execution environments), and uses zero-knowledge authentication to avoid revealing credentials, while also reporting performance measurements and formal verification tooling choices (Noise framework with ProVerif) \cite{r2025}. These two papers converge on the claim that future-proofing E2EE requires rethinking the handshake/key establishment layer while maintaining an application-level E2EE experience, but they diverge in emphasis: Sharma et al.\ foreground protocol-level replacement of vulnerable key exchange in Signal-like settings, whereas EnConvo couples PQ key exchange with broader endpoint hardening and authentication mechanisms as an integrated application stack \cite{sharma2023,r2025}.

A third direction interrogates how E2EE claims can be undermined at endpoints via client-side scanning (CSS). Alsop et al.\ argue that CSS effectively scans all message modalities on sending/receiving devices and therefore weakens the practical meaning of E2EE in mainstream messengers that advertise it \cite{alsop2025}. Their proposed response is architectural: offload cryptographic operations to an external, secure device (augmented reality glasses) to reduce client-side vulnerabilities and better preserve tamper-resistant private communication in sensitive contexts \cite{alsop2025}. This contributes a counterpoint to protocol-centric improvements: even if key exchange and cryptography are strengthened (including post-quantum), endpoint inspection can negate E2EE’s privacy guarantees, pushing research toward trusted hardware separation as a complement to cryptographic design \cite{alsop2025,r2025}.

A fourth research thread tackles a core governance dilemma: E2EE hampers detection and mitigation of abusive or malicious content because intermediaries cannot inspect ciphertext. Jiang et al.\ position existing accountability mechanisms (e.g., message franking/source tracing) as limited because they support only plain-data moderation under E2EE, and propose sMAC to enable sanitization and moderation over encrypted messages while aiming to preserve privacy, sender anonymity, and backward security \cite{jiang2024}. By instantiating sMAC with “amenable ACE” and providing formal proofs plus efficiency claims independent of message size for key operations, this work exemplifies an attempt to reconcile encrypted-content protection with moderated ecosystems without reverting to plaintext access \cite{jiang2024}. In contrast to CSS-based approaches (which inspect data at endpoints), sMAC targets cryptographic mechanisms that permit certain control and inspection properties without broadly exposing message contents, thereby articulating an alternative path for platforms seeking both E2EE and abuse resistance \cite{jiang2024,alsop2025}.

Finally, the literature reflects continued attention to the practical costs of E2EE. Thapa et al.\ compare resource consumption (CPU/GPU/memory) during encrypted messaging in WhatsApp, Telegram, and Signal, framing E2EE as imposing device-side computational demands and proposing CPU offloading to reduce resource utilization while maintaining confidentiality \cite{thapa2025}. This efficiency perspective complements the security-expansion directions above: stronger cryptography (e.g., post-quantum handshakes) and added protections (e.g., TEEs or external devices) may increase complexity or overhead, making performance analysis and optimization an enabling concern for deployability \cite{thapa2025,r2025,alsop2025}.

Across the set, the Signal protocol functions both as a reference point and as a motivation for evolution. Oppliger documents its central role and associated privacy issues (including metadata exposure), while later works either explicitly build “Signal-like” E2EE into new architectures, critique threats that bypass protocol guarantees (CSS), or propose cryptographic extensions and replacements (post-quantum key exchange; encrypted moderation primitives) that aim to preserve E2EE’s user-facing promise while adapting to new technical and societal constraints \cite{oppliger2020,kanimozhi2025,alsop2025,sharma2023,jiang2024}. Meanwhile, more introductory treatments still describe E2EE via public/private key pairs and the “signal encryption protocol” for secure data transmission, reflecting continued reliance on the conceptual framing of public-key-based E2EE even as specialized research pushes into post-quantum and accountability-aware designs \cite{p2023}.

\subsection{forensics}
Recent Signal-related instant messaging research in forensics has emphasized post-compromise evidence recovery on Android by targeting application-resident artifacts rather than attempting to weaken end-to-end encryption in transit \cite{krishnapriya2021}. Krishnapriya et al. frame Signal’s growing adoption as increasing the volume of potential digital evidence while simultaneously raising acquisition difficulty due to stronger privacy protections compared with other mainstream messengers \cite{krishnapriya2021}. In response, their contribution is methodological: they propose acquiring Signal data via decryption of the application’s backup file to recover chat messages and attachments from Android devices, shifting the forensic locus from network interception to device-level extraction \cite{krishnapriya2021}. This positioning implicitly treats backups as a pivotal evidentiary surface for Signal on Android, where successful forensic analysis depends on the ability to access and decrypt stored application data to reconstruct conversational content and associated media \cite{krishnapriya2021}.

\subsection{formal-methods}
Since 2020, formal-methods work around Signal-style secure messaging has diversified beyond conventional confidentiality proofs toward models that make the ``human in the loop'' and knowledge reasoning explicit. Two complementary directions are visible: (i) logic-based analyses that treat security properties as statements about what principals and adversaries can know, and (ii) formal game-based models that integrate user-mediated authentication with modern ratcheting key exchange and post-compromise settings \cite{kamkuemah2021,dowling2021}.

A key point of contrast is how authentication is represented and stress-tested. Kamkuemah frames authentication and secrecy for Signal using epistemic logic, emphasizing a design-first workflow where desired properties are expressed prior to implementation and then checked against protocol behavior \cite{kamkuemah2021}. In this knowledge-centric view, Signal is found secure with respect to the expressed authentication and secrecy properties, illustrating how epistemic formalisms can validate protocol goals in terms of ``who knows what'' \cite{kamkuemah2021}. Dowling and Hale, however, argue that authenticity in real deployments depends on user-mediated verification of long-term keys and that this layer has often been assumed away in formal analysis; they introduce a formal model capturing user-mediated entity authentication that composes with ratcheted key exchange and explicitly targets active, post-compromise man-in-the-middle threats \cite{dowling2021}. Under that model, Signal's user-mediated authentication cannot be proven secure, motivating a concrete protocol-level fix to enable detection of an active adversary \cite{dowling2021}. Taken together, these works suggest that the choice of formalism (epistemic property checking vs.\ composable authentication modeling) can determine whether Signal appears secure, particularly when adversarial interaction with user verification is treated as part of the protocol surface \cite{kamkuemah2021,dowling2021}.

In parallel, formal methods are also being used to extend secure-messaging functionality while preserving privacy and accountability under encryption. Jiang et al.\ propose sMAC via a new primitive (amenable ACE) to support sanitization and encrypted-message moderation, and provide formal security proofs in the standard model for privacy, sender anonymity, and backward security \cite{jiang2024}. Although not a Signal-protocol analysis per se, this line complements Signal-focused formalization by showing how standard-model proof techniques are being applied post-2020 to reconcile end-to-end encryption with moderation/accountability goals, expanding the formal-methods agenda from ``secure messaging as confidentiality'' to richer enforcement and abuse-handling objectives \cite{jiang2024}. In synthesis, the post-2020 formal-methods landscape couples deeper modeling of human-mediated authentication failure modes in Signal-like systems with rigorous proofs for emerging encrypted-messaging services that broaden the design space beyond traditional secrecy and forward/backward security guarantees \cite{dowling2021,jiang2024}.

\subsection{formal-modeling}
Since 2020, formal-modeling work around instant messaging and the Signal protocol spans both (i) descriptive/formal representations of protocol operation and (ii) proof-oriented models tied to new privacy properties and protocol variants. An automata-theoretic formalization is used to make the organization of Signal-based online call signaling explicit via nondeterministic event-driven automata, with the stated goal of clarifying protocol structure, relevant factors, and potential developer-facing problems \cite{vladimirovich2020}. In contrast, a more recent line of work leverages information-flow control techniques to construct formal models that support reasoning about metadata privacy for selected traffic classes, culminating in the design of DenIM as a Signal-variant enabling deniable instant messaging \cite{nelson2024}.

These efforts differ in what “formalization” is intended to achieve. The automata approach focuses on interpretability of the call-signaling organization—aiming to expose protocol organization and issues more clearly to practitioners by modeling behavior as an event-driven nondeterministic automaton \cite{vladimirovich2020}. The DenIM work uses formal modeling as a foundation for security arguments: it introduces a principled modeling approach for systems that provide metadata privacy for some, deniable, traffic and proves that such deniable traffic attains metadata privacy even against strong adversaries \cite{nelson2024}. This positions formalism not only as a documentation or design aid, but as a mechanism to connect system structure to a specific adversary model and property proof \cite{nelson2024}.

A second point of contrast is the level at which the model informs protocol evolution. The automata-based formalization is presented as a way to show protocol organization and potential problems, but it is not described as yielding new protocol mechanisms \cite{vladimirovich2020}. Conversely, the information-flow-based modeling directly supports protocol design: DenIM is explicitly framed as a variant of the Signal protocol enabling metadata privacy for deniable traffic, and the paper argues that state-of-the-art protocols can be extended to support this form of metadata privacy \cite{nelson2024}. Moreover, DenIM’s formal reasoning is coupled to an implementation strategy that preserves unmodified Signal traffic properties (e.g., low latency and existing features) while adding deniable-traffic support, tying the formal model to deployability constraints in real instant-messaging settings \cite{nelson2024}.

Taken together, these papers indicate that post-2020 formal-modeling related to Signal and instant messaging is bifurcating into (a) operational formalisms meant to render protocol behavior and developer-relevant pitfalls transparent and (b) security-driven formalisms that integrate with protocol modifications and proofs for properties beyond classic confidentiality, notably metadata privacy for subsets of traffic \cite{vladimirovich2020,nelson2024}. The latter also demonstrates an explicit bridging of information-flow control with anonymous-communication style guarantees, suggesting an expanded formal vocabulary for reasoning about messaging protocols in adversarial network environments \cite{nelson2024}.

\subsection{formal-verification}
Since 2020, formal verification work explicitly connected to Signal-derived messaging has begun to target settings beyond traditional one-to-one messengers, using symbolic-model tooling to validate security goals under asynchronous and multi-device constraints \cite{basem2023}. Basem et al.\ position Stick as a Signal-based protocol designed for social network platforms, and treat verification not as an afterthought but as a core step alongside protocol design and implementation \cite{basem2023}. Their use of Verifpal illustrates a trend toward mechanized symbolic analysis for rapidly iterating on protocol variations while retaining recognizably Signal-inspired properties such as authentication and confidentiality \cite{basem2023}.

Within this line of work, the distinctive verification focus is on whether Signal-like guarantees can be preserved when the communication model shifts from stable, synchronous sessions to many-to-many, asynchronous interaction across multiple devices \cite{basem2023}. Stick’s formally analyzed design introduces a new notion of session, multiple pairwise sessions, and refreshing identity keys, and verification is used to substantiate claims that sessions can be re-established while maintaining authentication and confidentiality \cite{basem2023}. In effect, the formal analysis is leveraged to argue that extending Signal-based constructions to social-network-style usage requires explicit protocol mechanisms for session recovery and key evolution, and that these mechanisms can be validated in a symbolic framework \cite{basem2023}.

A further synthesis point is how verification is tied to compromise recovery goals in contexts where group-oriented protocols often fall short \cite{basem2023}. Basem et al.\ argue, via symbolic verification, that Stick achieves forward secrecy and adds backward secrecy, supporting a form of post-compromise security in many-to-many communications \cite{basem2023}. This suggests an emerging verification agenda: rather than only proving baseline secrecy/authentication, formal methods are being applied to evaluate whether Signal-derived designs can support re-establishable sessions and compromise recovery properties in large-scale, social networking interaction patterns \cite{basem2023}.

\subsection{group-messaging}
Since 2020, group-messaging research around Signal’s ecosystem and related secure-messaging work has coalesced around three intersecting concerns: (i) how to extend two-party ratcheting ideas to dynamic groups, (ii) how to scale group security via alternative group keying architectures, and (iii) how to audit or analyze group systems whose encryption protects content but leaves room for integrity or metadata failures \cite{oppliger2024,ayed2025,sarvaiya2025,brigham2025,bobrysheva2021}. A common throughline is that “making Signal-like security work for groups” is treated as a distinct design space with different failure modes than one-on-one messaging, motivating both new protocol constructions and analyses of real deployments \cite{oppliger2024,ayed2025,brigham2025}.

One line of work explicitly generalizes ratcheting-based secure messaging to groups by adapting the Double Ratchet paradigm, but in markedly different directions. Bobrysheva and Zapechnikov argue that group chats require rethinking primitives under a post-quantum threat model, proposing a group key establishment approach based on isogenies of elliptic curves and framing the resulting construction as an “extended Double Ratchet” suitable for group communications \cite{bobrysheva2021}. In contrast, Ayed and Dawood pursue scalability within conventional (non-post-quantum) primitives by introducing a multiple ratchet group key agreement (MRGKA) that partitions a large group into subgroups (capped at 15 members) and treats each subgroup as a single entity for intra-group security management \cite{ayed2025}. Read together, these designs highlight a core tension: ratchet-inspired group protocols either prioritize future-proofing against quantum adversaries via new hardness assumptions \cite{bobrysheva2021} or prioritize operational scalability and flexible membership changes by restructuring group state and key distribution \cite{ayed2025}. Both works nevertheless converge on the premise that group messaging must preserve key security properties associated with the Signal-style ratcheting approach while tolerating asynchronous communication and frequent group churn \cite{bobrysheva2021,ayed2025}.

A complementary scaling narrative is articulated through the contrast between Signal and Messaging Layer Security (MLS). Oppliger positions Signal as fundamentally two-party in its original design, emphasizing that extending comparable guarantees to large groups is non-trivial and motivating MLS as a protocol designed for groups “for thousands of participants” \cite{oppliger2024}. This reframing helps contextualize why ratchet-based group extensions (including those above) coexist with MLS-oriented approaches: rather than a single “best” group design, the literature recognizes distinct architectural commitments—adapting two-party secure messaging patterns to groups versus adopting group-native key management aimed at large deployments \cite{oppliger2024,ayed2025,bobrysheva2021}.

Beyond key establishment and scaling, recent work stresses that group-messaging security also depends on properties that are not fully addressed by end-to-end encryption alone, such as message ordering integrity or metadata privacy. Sarvaiya and Burger examine cross-platform interoperable group messaging in the IETF MIMI setting, where MLS ciphertexts are routed through a central hub that timestamps and broadcasts messages \cite{sarvaiya2025}. They show that even if content confidentiality holds, a compromised hub can silently drop, delay, or reorder messages, and they propose a lightweight Merkle-tree-based audit mechanism that clients embed into encrypted application messages to detect hub misbehavior \cite{sarvaiya2025}. This audit-oriented perspective complements Oppliger’s framing of MLS as a scaling solution by highlighting that deploying group protocols in hub-mediated interoperability environments introduces additional trust and integrity risks beyond core group key management \cite{oppliger2024,sarvaiya2025}.

In parallel, analysis of Signal’s group features underscores that metadata can remain a critical vulnerability even when the protocol aims to hide group structure from the server. Brigham and Hopper target Signal’s Sealed Sender and Private Groups mechanisms and present an attack that links groups of communicating entities using recipient metadata alone, undermining the intended metadata-hiding guarantees in group conversations \cite{brigham2025}. Their finding connects to the broader theme that group-messaging security is multi-dimensional: strengthening cryptographic key evolution (as in ratchet extensions) does not automatically secure group privacy against traffic analysis, and server-mediated designs can still leak exploitable metadata \cite{brigham2025,oppliger2024}. Notably, the mitigation they discuss is “server-agnostic” and user-deployable, emphasizing practical tradeoffs between overhead and protection—an echo of the overhead/efficacy tradeoffs made explicit in audit layers for hub integrity in interoperable group messaging \cite{brigham2025,sarvaiya2025}.

Taken together, the post-2020 literature in this set suggests a maturing understanding of group messaging as an ecosystem problem spanning cryptographic design (including post-quantum directions), scalable group architectures, and operational attack surfaces such as hub behavior and metadata leakage \cite{bobrysheva2021,ayed2025,oppliger2024,sarvaiya2025,brigham2025}. The resulting research agenda is less about merely “adding groups” to Signal-style secure messaging and more about specifying which group properties are prioritized—scalability, post-quantum resilience, interoperability integrity, or metadata privacy—and then engineering mechanisms whose costs and assumptions align with those priorities \cite{oppliger2024,ayed2025,bobrysheva2021,sarvaiya2025,brigham2025}.

\subsection{instant-messaging}
Since 2020, instant-messaging research connected to the Signal protocol has increasingly treated the IM stack as an ecosystem in which confidentiality of message contents is only one property among many (e.g., availability, authentication, metadata privacy, and device integrity), leading to work that either (i) redesigns surrounding system components while keeping Signal-compatible primitives, or (ii) evaluates how real deployments leak or fail despite end-to-end encryption (E2EE) \cite{nelson2024,jain2021,bozorgi2023}. This shift is visible in studies that foreground non-cryptographic pressures of IM at scale (e.g., signaling overhead from always-online heartbeats) alongside traditional security goals, suggesting that “secure messaging” must also be operationally efficient to be deployable \cite{yi2020,thapa2025}.

A recurring theme is that Signal-style E2EE is necessary but insufficient once adversaries exploit endpoints, infrastructure control points, or side channels. On the endpoint side, vulnerability analysis of a widely used Signal-protocol messenger demonstrated that attacks can bypass protections to access confidential local user data (documents, media, audio memos), underscoring that application and OS-level exposure can nullify strong transport/content cryptography \cite{jain2021}. Complementing this, ARSecure explicitly targets client-side scanning (CSS) as a threat model in which messages are inspected at the sender/receiver despite E2EE claims by mainstream messengers (including Signal), proposing to move crypto operations onto an external AR device to reduce client compromise and tampering opportunities \cite{alsop2025}. Together, these works motivate a broadened notion of IM security in which “end-to-end” must include hardened clients or trusted peripherals, not only robust protocols \cite{jain2021,alsop2025}.

On the network side, multiple papers show that encryption does not prevent inference from traffic patterns. Large-scale traffic-analysis work across popular IM apps (including Signal) demonstrates that adversaries observing encrypted traffic can infer sensitive user activities and identify participants in target communications, proposing IMProxy as a client-side countermeasure without requiring provider support \cite{bozorgi2023}. Parallel to attack development, the field has also invested in tooling to make encrypted-traffic research reproducible: an automated Android-based traffic generation and labeling system produces datasets for several IMAs including Signal, enabling data-driven encrypted traffic analysis where public datasets are scarce \cite{erdenebaatar2023}. These efforts jointly indicate maturation from isolated demonstrations to an empirical pipeline (data generation $\rightarrow$ attack evaluation $\rightarrow$ countermeasure testing), while keeping the focus on IMAs’ observable network behavior rather than message contents \cite{erdenebaatar2023,bozorgi2023}.

A second cluster of work reframes what privacy should mean for IM beyond content secrecy, particularly metadata privacy and deniability. Nelson et al.\ propose a formal modeling approach that bridges information flow control with anonymous communication by introducing “deniable traffic,” then instantiate this idea as DenIM, a variant of the Signal protocol that can run on top of unmodified Signal to add metadata privacy for some traffic while preserving low latency for ordinary traffic and existing features \cite{nelson2024}. In contrast to obfuscation add-ons such as IMProxy that aim to reduce traffic-analysis signal without changing the underlying messenger, DenIM emphasizes protocol-level objectives and formal reasoning, while still targeting backward compatibility by layering over existing Signal deployments \cite{nelson2024,bozorgi2023}. This tension—deployable overlays versus protocol modifications—highlights an emerging design space where incremental adoption is treated as a first-class constraint for privacy-enhancing IM \cite{nelson2024}.

Authentication and key management have also been repositioned as critical weak points even when Signal solves message encryption. Morales et al.\ argue that while protocols like Signal address encryption, public-key authentication in IM remains problematic due to centralized infrastructure and user misunderstanding of E2EE; they propose removing IM providers as identification/authentication authorities by using Trust over IP with decentralized identifiers and verifiable credentials to manage public keys outside provider control \cite{morales2024}. This direction aligns with decentralization-oriented IM system work that critiques central servers as single points of failure and control: Mashru et al.\ propose a decentralized, cloud-hosted IM architecture with E2EE intended to be fault tolerant and reduce corporate dominance \cite{mashru2023}. Although these works intervene at different layers (identity infrastructure vs.\ message routing/hosting), they converge on the idea that IM trust should be distributed across verifiable systems rather than concentrated in providers, implicitly complementing Signal’s cryptographic guarantees with governance and authentication decentralization \cite{morales2024,mashru2023}.

Interoperability and group messaging have additionally introduced new integrity risks and protocol proposals adjacent to Signal. Sarvaiya and Burger analyze the IETF MIMI interoperability setting where a central hub relays MLS ciphertexts for cross-platform group messaging (including services like Signal) and show that while E2EE may hold, the hub can still undermine order integrity by dropping, delaying, or reordering messages; they propose a Merkle-tree audit layer implemented on OpenMLS to detect hub misbehavior efficiently \cite{sarvaiya2025}. In a related but distinct vein, Ayed and Dawood propose a multiple ratchet group key agreement protocol inspired by Double Ratchet principles to handle dynamic group membership by subdividing users into capped subgroups with additional authentication codes, emphasizing scalability and asynchronous operation \cite{ayed2025}. Taken together, these works show how the move from single-app, single-room assumptions to interoperable and large-group contexts shifts security priorities toward auditable delivery and scalable key evolution, rather than only per-message secrecy \cite{sarvaiya2025,ayed2025}.

At the implementation and performance level, some post-2020 IM research continues to adapt Signal components (e.g., X3DH and Double Ratchet) while proposing alternative key exchange or acceleration strategies, reflecting pressure to balance confidentiality with resource constraints. One line proposes integrating a new key exchange (ORDEX) with X3DH and Double Ratchet and using AES for message encryption, claiming substantial speedups over ECDH-based approaches as a way to reduce security-induced latency \cite{m2024}. Another evaluates resource consumption (CPU/GPU/memory) of encrypted messaging in WhatsApp, Telegram, and Signal and proposes CPU offloading to reduce device burden while maintaining “acceptable” confidentiality \cite{thapa2025}. These performance-driven proposals sit in productive tension with earlier observations that IM’s always-online design already imposes energy and signaling costs on mobile networks (e.g., heartbeat messaging), motivating system-level relaying and aggregation approaches like eDirect to cut signaling traffic and handset energy use \cite{yi2020,thapa2025}. Across these works, the implicit synthesis is that cryptographic robustness alone does not guarantee an acceptable IM user experience or deployability; efficiency improvements are treated as security-enabling because they reduce incentives to disable protections or avoid secure platforms \cite{thapa2025,yi2020}.

Finally, several papers broaden the instant-messaging security landscape beyond Signal-centric design by analyzing adjacent messaging infrastructures and application domains that nevertheless intersect with IM security assumptions. Zhao et al.\ provide an in-depth security study of 5G RCS, finding weak cellular ID binding that enables persistent hijacking and makes impersonation/eavesdropping feasible even when E2EE is present, illustrating how identity binding and carrier infrastructure can dominate outcomes regardless of cryptographic channel protections \cite{zhao2022}. Samper and Ferreira’s systematization of “safer sexting” across messaging and social apps further emphasizes that IM privacy harms often arise from downstream sharing, accountability gaps, and feature design, not only encryption correctness, reinforcing the broader trend toward socio-technical threat models for private chat use \cite{samper2025}. Within this widened scope, Signal protocol research since 2020 appears less about the core cipher suite itself and more about making IM secure in practice by hardening clients, limiting metadata leakage, decentralizing authentication authority, and preserving integrity and usability under new deployment models such as interoperability and large groups \cite{nelson2024,morales2024,sarvaiya2025}.

\subsection{interoperability}
Work since 2020 on interoperability in end-to-end encrypted instant messaging highlights the tension between cross-platform connectivity and the new trust assumptions introduced by interoperability infrastructure. In the IETF More Instant Messaging Interoperability (MIMI) setting, a central hub mediates group messaging across otherwise siloed services (e.g., WhatsApp, Signal, Telegram) by timestamping and broadcasting Messaging Layer Security (MLS) ciphertexts, preserving payload confidentiality while making ordering and delivery properties dependent on hub behavior \cite{sarvaiya2025}. Sarvaiya and Burger identify this hub as a single point where compromise can silently induce message drops, delays, or reordering without breaking end-to-end encryption, reframing interoperability risk from cryptographic compromise to integrity failures in the delivery substrate \cite{sarvaiya2025}.

To address this interoperability-specific threat model, \cite{sarvaiya2025} proposes an audit layer that is explicitly designed to be lightweight and client-driven rather than relying on additional trusted infrastructure, which aligns with the practical constraints of deploying across heterogeneous providers participating in MIMI. Their approach treats the hub-assigned timestamp stream as an auditable log: each client maintains an ordered list of received messages with timestamps and periodically generates Merkle-tree proofs over this local view, disseminating proofs by embedding them inside encrypted application messages \cite{sarvaiya2025}. This design choice leverages the existing end-to-end channel (MLS-protected application messages) to propagate accountability information, aiming to keep interoperability compatible with existing security boundaries while adding a verifiable mechanism for detecting hub misbehavior \cite{sarvaiya2025}.

A key synthesis point is how \cite{sarvaiya2025} counters a malicious hub that might selectively suppress evidence of its own misconduct. By having clients probabilistically sample messages to carry proofs, and by ensuring the hub cannot predict which ciphertexts contain audits, the scheme is intended to make selective dropping of audit-carrying messages difficult without broadly degrading service \cite{sarvaiya2025}. The protocol further emphasizes collective verification: upon receipt of a proof, other clients verify and emit their own proofs, so inconsistencies propagate to the full group and create a non-repudiable record of hub misconduct within the room \cite{sarvaiya2025}. In this way, interoperability is treated not merely as a routing problem but as a distributed accountability problem where correctness must be reconstructed from multiple client perspectives rather than assumed from the interoperability hub \cite{sarvaiya2025}.

Finally, the evaluation reported in \cite{sarvaiya2025} suggests that such accountability can be added without undermining deployability: a Rust prototype on OpenMLS in a 100-node emulated network shows high detection rates under probabilistic attacks and low overhead in both memory and processing for proofs \cite{sarvaiya2025}. Taken together, this work indicates that post-2020 interoperability research is beginning to operationalize integrity and auditability as first-class requirements alongside end-to-end encryption, particularly when central hubs are introduced to bridge previously closed messaging ecosystems \cite{sarvaiya2025}.

\subsection{literature-review}
Since 2020, literature-review work on instant messaging and the Signal protocol has increasingly served a dual role: consolidating technical understanding of established E2EE mechanisms while mapping emerging pressures that challenge their long-term adequacy, notably scalability to groups and post-quantum risk \cite{oppliger2024,boras2025}. Across the provided works, the common thread is that Signal is treated as a baseline (or reference point) for modern secure messaging, but the review literature diverges in what it considers the primary frontier: protocol engineering and standardization (Signal-to-MLS), cryptographic migration (classical-to-post-quantum), or socio-technical safety features around sensitive content shared over private chats \cite{oppliger2024,boras2025,samper2025}.

Oppliger’s book-length treatment positions the Signal protocol as widely deployed yet underexplored in mainstream security literature, and frames its contribution as a structured explanation meant to bridge that gap \cite{oppliger2024}. In doing so, it links the review of Signal’s two-party design to the problem of extending E2EE guarantees to large groups, treating group messaging not merely as an implementation detail but as a distinct design challenge \cite{oppliger2024}. This emphasis complements Boras’ literature review, which assumes Signal’s centrality in secure messaging and instead foregrounds a different looming design constraint: the anticipated fragility of classical public-key primitives under practical quantum computation and the resulting need for quantum-resistant primitives that could be integrated into secure instant messaging protocols such as Signal \cite{boras2025}. Taken together, these two reviews portray a research landscape in which “Signal-protocol-related” work after 2020 is not only about analyzing a fixed protocol, but also about identifying migration paths—either toward protocols designed for group-scale settings (MLS) or toward post-quantum cryptographic building blocks—while maintaining the security expectations associated with Signal-style E2EE \cite{oppliger2024,boras2025}.

Whereas Oppliger and Boras are anchored in protocol and cryptographic design, Samper and Ferreira’s SoK broadens the literature-review lens from protocol correctness and adversarial cryptography to the end-to-end safety and privacy outcomes of instant messaging use cases, specifically sexting and self-generated nudes \cite{samper2025}. Their systematic literature review (10{,}026 search results plus cross-references) and accompanying survey of OS features and 52 real-world apps highlights a diverse ecosystem of proposals and deployed features intended to provide security, privacy, or accountability benefits when sharing sensitive media \cite{samper2025}. This body of work implicitly complements Signal-centric protocol literature by emphasizing that private-chat security is not exhausted by transport-level confidentiality: the review organizes a sexting threat model and a taxonomy of mechanisms, and argues that no single technique fully addresses the full threat space \cite{samper2025}. In synthesis with Signal-focused reviews, this indicates a post-2020 trend in which literature reviews increasingly treat “secure messaging” as a stack that spans cryptographic protocols (Signal), scalable group messaging designs (MLS), and application/OS-level features and threat models that shape user-facing privacy and harm \cite{oppliger2024,samper2025}. Meanwhile, the post-quantum review agenda adds another cross-cutting constraint: even if messaging features and group protocols improve, their long-term guarantees still depend on swapping out vulnerable classical primitives for quantum-resistant alternatives within protocols like Signal \cite{boras2025}.

Overall, the reviewed literature suggests that post-2020 scholarship has shifted from treating Signal primarily as a single protocol to treating it as a reference design embedded in a broader evolution of secure instant messaging: reviews codify its mechanisms for practitioners, juxtapose it with MLS for large-group scalability, and motivate cryptographic modernization for quantum resistance, while parallel SoK work maps broader privacy and safety interventions in private-chat ecosystems \cite{oppliger2024,boras2025,samper2025}.

\subsection{metadata-privacy}
Recent work since 2020 on metadata privacy in Signal-like instant messaging highlights a tension between (i) protocol mechanisms intended to conceal communication relationships and (ii) traffic-analysis vectors that can re-identify those relationships from residual metadata. One line of research reframes the design objective to make metadata privacy achievable for only a subset of traffic, enabling pragmatic deployment without requiring universal, high-overhead transport-layer anonymity \cite{nelson2024}. A complementary line evaluates deployed metadata-hiding features in Signal and demonstrates that even when sender identity and group membership are ostensibly concealed, recipient-side metadata patterns can still link participants and groups \cite{brigham2025}.

Nelson et al.\ argue that adopting strong transport-layer privacy for all traffic faces adoption barriers (notably mobile performance costs), motivating a new design space in which only some traffic is protected while other traffic remains “deniable” with respect to metadata \cite{nelson2024}. Their contribution is not only a protocol sketch but a principled modeling approach that bridges information-flow control with anonymous communication: they formalize conditions under which deniable traffic attains metadata privacy against strong adversaries \cite{nelson2024}. In this framework, Signal is treated as an extensible substrate: they design DenIM as a Signal-variant and show in a proof-of-concept implementation that DenIM can run on top of unmodified Signal while preserving low latency for regular Signal traffic and adding support for deniable (metadata-private) traffic \cite{nelson2024}. This “selective” privacy objective thus trades universal coverage for deployability and performance, and treats coexistence with standard messaging features as a first-class requirement \cite{nelson2024}.

Brigham and Hopper, by contrast, interrogate the extent to which Signal’s existing metadata-hiding extensions (Sealed Sender and Private Groups) meet their intended goals when confronted with traffic analysis focused on group conversations \cite{brigham2025}. They show that groups of communicating entities can be linked through recipient metadata alone, thereby defeating the stated goal of concealing which parties communicate in groups from compromised servers \cite{brigham2025}. Their analysis further indicates that a prior defense aimed at a similar linking attack for Sealed Sender pairs does not carry over to the group setting, underscoring that metadata protections may be fragile under changes in communication structure (pairwise vs.\ group messaging) \cite{brigham2025}. They then outline a user-implementable, server-agnostic mitigation that explicitly illustrates a cost curve: increased communication overhead can improve resistance to linking, but no “free” defense is implied \cite{brigham2025}.

Taken together, these papers suggest that metadata privacy in instant messaging is increasingly approached as an end-to-end systems property rather than a single mechanism: formal privacy goals and selective deployment strategies aim to make protection practical \cite{nelson2024}, while empirical/theoretical attacks reveal that real-world features can leak sufficient structure (e.g., recipient metadata patterns) to reconstruct sensitive relationships even when sender and group identifiers are nominally hidden \cite{brigham2025}. This juxtaposition highlights a core design challenge for Signal-derived protocols: extending the protocol to add metadata-private modes requires not only provable properties in a model but also careful consideration of what metadata remains observable (and exploitable) in concrete deployments \cite{nelson2024,brigham2025}.

\subsection{multi-device}
Recent Signal-protocol-related work on multi-device settings has moved beyond “device cloning” assumptions in messengers toward supporting asynchronous, re-establishable sessions across multiple endpoints, motivated by social-network interaction patterns \cite{basem2023}. Basem et al. frame multi-device not merely as key distribution to additional devices, but as a need to repeatedly recover secure communication state when devices join, go offline, or are replaced, while still retaining Signal-derived security properties \cite{basem2023}. Their Stick protocol operationalizes this by introducing a new session concept designed to be re-established, using multiple pairwise sessions (rather than a single monolithic session) as a structural mechanism for multi-device communications \cite{basem2023}.

A key synthesis point is that Stick treats multi-device support as intertwined with compromise recovery rather than as a usability add-on: it refreshes identity keys and aims to preserve forward secrecy while additionally introducing backward secrecy in an asynchronous multi-device setting \cite{basem2023}. This emphasis positions multi-device capability as a prerequisite for post-compromise security in many-to-many communications, and Stick claims to achieve “a form” of post-compromise security where many group protocols fall short \cite{basem2023}. Methodologically, the paper anchors these claims with symbolic-model formal verification in Verifpal and backs the feasibility of its multi-device design via an implemented open-source API and performance/usability evaluation in a social-network context \cite{basem2023}.

\subsection{post-quantum}
Across the post-2020 literature on instant messaging security, post-quantum work clusters around (i) adapting Signal’s core key-establishment and ratcheting ideas to quantum-resistant primitives and (ii) building new end-to-end encrypted systems that combine PQC with adjacent mechanisms (e.g., hardware protection and novel authentication), often benchmarking performance and verification rather than proving composable equivalence to Signal’s security goals \cite{boras2025,angom2025,r2025,sharma2023,bobrysheva2021}. A recurring premise is that the principal pressure point is the public-key component in widely deployed secure messaging (notably Signal’s Curve25519-based ECDH in X3DH-style setup), motivating hybrid or replacement designs that aim to preserve asynchronous messaging and forward/future secrecy under quantum adversaries \cite{sharma2023,boras2025}.

One line of work targets Signal-style initial key establishment. Angom et al.\ explicitly analyze Signal’s post-quantum extended Diffie--Hellman (PQXDH) framing and propose MLXDH, integrating NIST-standardized primitives ML-KEM and ML-DSA into the initial exchange as a proof-of-concept hardening step \cite{angom2025}. Their motivation highlights a nuanced gap: even when a post-quantum mechanism is introduced for key agreement, mutual authentication can remain tied to discrete-logarithm assumptions, leaving a partially classical dependency that they seek to further harden via the addition of PQ signatures \cite{angom2025}. In contrast, Sharma et al.\ frame the core need more directly as replacing ECDH key exchange in Signal-protocol-based applications with a quantum-safe alternative, positioning PQC substitution at the heart of maintaining Signal’s broader security properties (e.g., forward secrecy, deniability, asynchronous operation), though without the same emphasis on the authentication subproblem singled out by MLXDH \cite{sharma2023}. Read together, these papers suggest an emerging design tension: incremental hardening via hybridization and PQ add-ons can be engineered and evaluated quickly, but achieving a \emph{fully} post-quantum Signal-style handshake requires resolving authentication (and not only key agreement) under PQ assumptions \cite{angom2025,sharma2023}.

A second line extends post-quantum thinking beyond 1:1 setup into group messaging, where maintaining ratchet-style security properties becomes more complex. Bobrysheva and Zapechnikov propose a post-quantum secure group key establishment approach based on elliptic-curve isogenies, described as an extension of the Double Ratchet protocol to support group chats \cite{bobrysheva2021}. This proposal complements the handshake-focused efforts by emphasizing that post-quantum migration is not only a matter of swapping primitives in X3DH/PQXDH-like setup, but also of rethinking key management in multi-party settings while preserving the expected security properties of modern messengers \cite{bobrysheva2021}. In this sense, group messaging work broadens the threat-driven agenda (quantum attacks) into a functionality-driven agenda (group chat construction), whereas the Signal key-establishment hardening work stays closer to the protocol’s initial trust and authentication bootstrapping \cite{bobrysheva2021,angom2025}.

A third direction is application/system design that treats PQC as one component in a larger secure messaging stack. EnConvo combines Kyber-1024 for key exchange (explicitly retaining X25519 alongside Kyber for backward compatibility), HKDF-SHA3-512 for key derivation, and hardware-based AES-256 in OCB mode for bulk encryption, while adding hardware enclaves (Intel SGX/ARM TrustZone) for key protection and zk-SNARK-based passwordless authentication \cite{r2025}. The emphasis here is not on preserving Signal-protocol compatibility per se, but on assembling a PQ-oriented end-to-end encrypted messenger with performance measurements (e.g., Kyber exchange latency and per-message encryption time) and formal verification of protocol execution using the Noise framework with ProVerif \cite{r2025}. Compared with MLXDH and the Double-Ratchet extension, EnConvo exemplifies a more systems-oriented response to the quantum threat model, where migration risk is handled through hybrid compatibility (Kyber with X25519) and assurance is sought through tooling and deployment-oriented security controls (enclaves), rather than by deriving security directly from Signal’s established protocol proofs or threat model alignment \cite{r2025,angom2025,bobrysheva2021}.

Finally, Boras’ literature review provides a unifying lens by framing these efforts as part of broader attempts to develop, evaluate, and optimize quantum-resistant primitives for secure instant messaging protocols, with particular attention to Signal \cite{boras2025}. Within the set of papers considered here, that framing is reflected in (a) targeted protocol constructions for Signal’s initial exchange using NIST PQ primitives \cite{angom2025}, (b) ratchet-inspired group extensions using alternative PQ assumptions (isogenies) \cite{bobrysheva2021}, and (c) full application frameworks that combine PQC with verification and hardware mechanisms while preserving interoperability through hybrids \cite{r2025}. Collectively, post-2020 work thus converges on the need for post-quantum upgrades but diverges on where to intervene (handshake vs.\ group state evolution vs.\ system stack) and on whether “post-quantum Signal” is best pursued via protocol-preserving substitutions/hardening or via new messenger designs that borrow only selected architectural elements \cite{boras2025,angom2025,bobrysheva2021,r2025,sharma2023}.

\subsection{privacy-leakage}
Recent work since 2020 shows privacy leakage in secure messaging arising from two complementary sources: \emph{metadata/traffic} that remains observable despite end-to-end encryption, and \emph{functional extensions} (e.g., moderation/accountability) that must carefully balance privacy with governance requirements. On the first axis, Bozorgi et al.\ demonstrate that encrypted IM traffic itself can be sufficient to infer sensitive user activities: their traffic-analysis attacks identify participants of target communications (including forums) with high accuracy while assuming no app vulnerabilities, underscoring that confidentiality of message content does not imply confidentiality of interaction patterns \cite{bozorgi2023}. On the second axis, Jiang et al.\ argue that moderation and sanitization are difficult over ciphertext and that existing approaches (e.g., franking/source tracing) only support moderation over plain data in end-to-end encrypted settings, motivating cryptographic redesigns that explicitly incorporate accountability while aiming to preserve privacy properties \cite{jiang2024}.

Across these directions, a common synthesis point is that “privacy leakage” is increasingly framed as a system-level property rather than a pure cryptographic one. Bozorgi et al.\ operationalize leakage as an adversary’s ability to infer who is communicating in which contexts from encrypted network traces, positioning the network observer as a realistic threat actor in surveillance/censorship environments \cite{bozorgi2023}. In contrast, Jiang et al.\ treat leakage as the risk introduced when platforms need to inspect or control abusive content flows: their sMAC framework is designed to enable moderation/sanitization \emph{over encrypted messages} while claiming data privacy, sender anonymity, and backward security, thereby shifting the locus of leakage from passive observation to the design of authorized inspection capabilities \cite{jiang2024}. Read together, these works suggest that privacy risks persist both when systems expose too much \emph{unintended} information (traffic patterns) and when they intentionally add \emph{intended} mechanisms (moderation/accountability) that, if poorly designed, could become new leakage channels \cite{bozorgi2023,jiang2024}.

The papers also contrast in how they mitigate leakage. Bozorgi et al.\ emphasize that standard countermeasures can degrade traffic-analysis effectiveness and propose an add-on client-side countermeasure (IMProxy) that requires no provider support, implying a pragmatic posture where leakage is mitigated through traffic obfuscation at the edge when protocols/apps remain unchanged \cite{bozorgi2023}. Jiang et al., conversely, pursue a protocol/primitive-level solution—amenable ACE within sMAC—supported by formal proofs and efficiency claims (costs independent of message size for key algorithms), reflecting a design-oriented approach where leakage is constrained by construction through carefully scoped authorization and inspection \cite{jiang2024}. This divergence highlights a broader trade-off in the post-2020 landscape: retrofitting defenses to protect users against network-side inference versus redesigning cryptographic access-control/moderation to prevent privacy erosion when adding governance features \cite{bozorgi2023,jiang2024}.

Finally, Samper and Ferreira’s systematization of “safer sexting” situates privacy leakage within a wider set of threats and mitigations observed in both research proposals and real-world apps, emphasizing that no single technique covers all threats and that privacy/accountability features must be composed with care \cite{samper2025}. Their SoK reinforces the synthesis above by framing private-chat risks as multifaceted—extending beyond content confidentiality to include broader privacy-related concepts and shortcomings of deployed features—thereby contextualizing why both traffic-analysis leakage and moderation-induced leakage remain central research concerns in modern messaging \cite{samper2025}.

\subsection{protocol-design}
Since 2020, protocol-design work around Signal and instant messaging has concentrated on (i) extending Signal-derived constructions beyond the baseline one-to-one setting, (ii) hardening the protocol stack against new threat models (notably post-quantum adversaries and metadata observers), and (iii) adding new functionality such as moderation and access control while retaining end-to-end privacy guarantees. Across these directions, a recurring design pattern is to preserve Signal’s core ratcheting and asynchronous usability while selectively replacing or augmenting components (e.g., initial key establishment, group key management, or the surrounding transport and governance mechanisms) to meet additional goals \cite{basem2023,nelson2024,angom2025,ayed2025,bobrysheva2021,jiang2024}.

A major cluster of designs targets the tension between Signal’s predominantly pairwise foundations and the realities of multi-party and multi-device communication. Stick explicitly positions itself as a Signal-based protocol tailored to social network platforms, emphasizing an asynchronous multi-device setting and introducing a new session concept with multiple pairwise sessions and identity-key refreshing to enable re-establishable sessions \cite{basem2023}. In contrast, group-focused designs attempt to lift double-ratchet-style security properties into scalable group key agreement structures. Bobrysheva and Zapechnikov propose an ``extended Double Ratchet'' for group key establishment and frame the design around post-quantum primitives based on isogenies \cite{bobrysheva2021}. Ayed and Dawood similarly keep Double Ratchet principles in view but pursue scalability through structural decomposition: users are partitioned into bounded subgroups that behave like single entities for intra-group communication, complemented with subgroup identity keys and signed pre-shared keys as protocol objects \cite{ayed2025}. Taken together, these works highlight two divergent responses to group complexity: Stick leans on many pairwise relationships and session re-establishment to cope with many-to-many interactions \cite{basem2023}, whereas MRGKA and the isogeny-based proposal restructure the group itself (via subgroups or group key establishment extensions) to manage dynamics while aiming to retain ratcheting-derived guarantees \cite{ayed2025,bobrysheva2021}.

Post-quantum hardening is pursued at multiple layers, but with different trade-offs regarding what parts of Signal are replaced versus augmented. Angom et al.\ focus narrowly on Signal’s initial key establishment, proposing MLXDH as a proof-of-concept construction that incorporates NIST’s ML-KEM and ML-DSA to increase cryptographic hardness, while acknowledging that mutual authentication in their design remains discrete-logarithm based and discussing paths toward full post-quantum authentication \cite{angom2025}. By comparison, Bobrysheva and Zapechnikov argue for post-quantum group messaging by specifying primitives and mechanisms for group key establishment based on isogenies, explicitly tying their approach to an extension of Double Ratchet rather than only modifying the initial handshake \cite{bobrysheva2021}. Sharma et al.\ motivate replacing ECDH key exchange (as used in Signal-like systems) with quantum-safe key exchange, reflecting a broader push to swap vulnerable public-key components while preserving Signal’s higher-level properties (e.g., forward secrecy and asynchronous operation) \cite{sharma2023}. Synthesizing these, the literature suggests a layered migration strategy: some proposals concentrate on post-quantum upgrades to the handshake (minimizing disruption to the rest of the protocol) \cite{angom2025}, while others more directly redesign group key establishment and ratcheting flows around post-quantum assumptions \cite{bobrysheva2021}, and still others frame the need to replace ECDH more generally within Signal-style architectures \cite{sharma2023}.

Another emerging protocol-design direction targets privacy leakage beyond message content, particularly metadata exposure. Nelson et al.\ argue that transport-layer data can leak who-communicates-with-whom metadata and propose a principled modeling approach that blends information flow control with anonymous communication. Their concrete outcome is DenIM, a Signal variant that supports metadata privacy for ``deniable'' traffic while maintaining low-latency and preserving unmodified Signal traffic and features when run atop unmodified Signal \cite{nelson2024}. This work contrasts with the cryptographic-hardening papers by treating metadata privacy as a system objective that reshapes design requirements, rather than as a property achieved solely by swapping primitives inside Signal’s cryptographic core \cite{nelson2024,angom2025}. It also complements group and multi-device designs by illustrating how protocol variants may coexist with baseline Signal operation, suggesting a pathway where privacy-enhancing traffic classes can be layered without abandoning the deployed protocol ecosystem \cite{nelson2024}.

Finally, protocol-design since 2020 has also expanded to include governance and safety functionality that traditionally sits outside Signal’s end-to-end encryption goals. Jiang et al.\ propose sMAC, a ``purified access control'' framework enabling sanitization and moderation over encrypted messages via an ``amenable ACE'' primitive, with formal security proof and efficiency claims that decryption, stamping, verification, and inspection costs are independent of message size \cite{jiang2024}. In synthesis with DenIM, this indicates two orthogonal expansions of the messaging protocol problem: one aims to reduce or deny metadata leakage during communication \cite{nelson2024}, while the other introduces encrypted-message moderation and accountability mechanisms without reverting to plaintext moderation \cite{jiang2024}. Both approaches broaden protocol design beyond classic confidentiality/integrity toward explicitly specified system-level objectives (metadata privacy or moderation under encryption), supported by formal reasoning or proofs as part of the design methodology \cite{nelson2024,jiang2024}.

Across these strands, the post-2020 literature portrays Signal not as a fixed protocol but as a design substrate: researchers either extend its ratcheting logic to groups and re-establishable multi-device sessions \cite{basem2023,ayed2025,bobrysheva2021}, retrofit it with post-quantum components at the handshake or broader key-establishment level \cite{angom2025,bobrysheva2021,sharma2023}, or wrap it with system mechanisms that address metadata privacy and encrypted moderation \cite{nelson2024,jiang2024}. The unifying protocol-design challenge is maintaining Signal’s deployed usability properties (asynchrony, feature compatibility, and performance) while altering foundational assumptions about adversaries, scale, and platform requirements \cite{basem2023,nelson2024,ayed2025,angom2025}.

\subsection{public-key-infrastructure}
Recent work highlights that, despite mature end-to-end encryption designs such as Signal, the remaining weak point in instant messaging security is public-key authentication, which is commonly mediated by centralized provider-controlled infrastructure \cite{morales2024}. In this framing, the public-key infrastructure problem is not primarily cryptographic; rather, it is rooted in the provider’s dual role as both service operator and de facto identification/authentication authority, creating trust dependencies and potential user confusion about what E2EE does and does not guarantee \cite{morales2024}. Morales et al.\ treat this as a structural mismatch between the security expectations of E2EE messaging and the governance model of key distribution, where provider involvement becomes a single locus for both trust and failure \cite{morales2024}.

To address this, the proposed direction is to externalize public-key authentication from the IM provider by replacing provider-centric key directories with a decentralized framework grounded in the Trust over IP (ToIP) model \cite{morales2024}. Specifically, Decentralized Identifiers (DIDs) and Verifiable Credentials (VCs) are used to manage and authenticate users’ public keys outside of the messaging platform’s administrative control, thereby removing the provider as an ``artificial'' authority for identification and authentication \cite{morales2024}. This shifts PKI-like functions (key binding and assurance) from a centralized, opaque mechanism to a decentralized and more transparent trust fabric, with the stated goal of improving user privacy, security, and autonomy in IM systems \cite{morales2024}. The work further positions decentralized key authentication as not only a technical redesign but also an interoperability and compliance enabler, arguing alignment with regulatory frameworks such as the Digital Markets Act (DMA) and eIDAS \cite{morales2024}.

\subsection{secrecy}
Since 2020, work on secrecy in Signal-style instant messaging has increasingly treated secrecy as intertwined with explicit models of knowledge and with the authentication mechanisms that gate who can hold which secrets. Kamkuemah reframes Signal’s secrecy analysis in a ``who knows what'' setting, using epistemic logic to express secrecy properties and then analyze protocol behavior against those properties, concluding that the protocol satisfies the desired secrecy guarantees when specified in this knowledge-centric manner \cite{kamkuemah2021}. This emphasis on expressing secrecy properties prior to implementation positions secrecy not only as a cryptographic outcome but as a formally stated constraint on adversarial and participant knowledge \cite{kamkuemah2021}.

In contrast, Dowling and Hale argue that secrecy guarantees in modern ratcheted messaging can be undermined when authenticity is compromised, particularly under active post-compromise conditions, because user-mediated key verification is typically modeled informally or assumed away \cite{dowling2021}. Their formal model targets user-mediated entity authentication composable with ratcheted key exchange and explicitly captures active post-compromise entity authentication security, highlighting that authenticity failures enable man-in-the-middle scenarios that can negate the practical meaning of secrecy for end users even when confidentiality mechanisms (e.g., ratcheting properties such as post-compromise security) are present \cite{dowling2021}. Within this model, they show Signal’s user-mediated authentication cannot be proven secure and propose a fix intended to allow detection of an active adversary, thereby restoring the conditions under which secrecy claims are meaningful in practice \cite{dowling2021}.

Together, these papers illustrate a methodological split in post-2020 secrecy research for Signal-related protocols: one line strengthens secrecy claims by refining the formal language used to state and verify secrecy (epistemic ``knowledge'' properties) \cite{kamkuemah2021}, while the other challenges secrecy as an end-to-end property unless authentication---including the human-in-the-loop verification channel---is formally captured under active compromise and proven to support the assumptions confidentiality proofs rely on \cite{dowling2021}. The resulting synthesis is that secrecy analyses are becoming less siloed: establishing secrecy increasingly depends on formally connecting what principals and attackers can know to whether the system can rule out (or at least detect) active impersonation paths that would otherwise redirect secrets to the wrong entity \cite{kamkuemah2021,dowling2021}.

\subsection{secure-messaging}

Since 2020, secure-messaging research around the Signal protocol has diversified from explaining and formalizing its core mechanisms to probing under-specified security properties (notably user-mediated authentication), extending ratcheting ideas to groups, and anticipating cryptographic transition pressures from post-quantum adversaries \cite{oppliger2020,vladimirovich2020,dowling2021,oppliger2024,sharma2023,boras2025,ayed2025}. Across this work, a recurring pattern is that confidentiality-centric guarantees (e.g., forward secrecy and post-compromise recovery) are treated as necessary but insufficient for real deployments, where usability constraints, group dynamics, and platform governance needs introduce new attack and design surfaces \cite{dowling2021,oppliger2024,jiang2024,samper2025}.

A first strand consolidates and operationalizes what “secure messaging” entails in practice by situating Signal within the broader E2EE ecosystem and highlighting adjacent risks such as metadata exposure and traffic analysis \cite{oppliger2020}. Complementing this systems-oriented framing, formalization efforts attempt to make protocol behavior more tractable for developers and analysts; for example, modeling Signal-based call setup with non-deterministic event-driven automata emphasizes protocol organization and potential implementation pitfalls, reflecting an ongoing push to bridge conceptual protocol descriptions and executable engineering concerns \cite{vladimirovich2020}. Oppliger’s later treatment further reinforces this role for integrative accounts by explicitly pairing Signal with Messaging Layer Security (MLS) as the emerging standardization pathway for scalable group E2EE, underscoring that secure-messaging research is increasingly about composable protocol stacks rather than a single two-party design \cite{oppliger2024}.

A second strand shifts attention from secrecy to authenticity in the presence of active attackers, interrogating a long-standing assumption in deployed secure messengers: that user-mediated key verification “solves” entity authentication \cite{dowling2021}. Dowling and Hale provide a formal model for user-mediated authentication that composes with ratcheted key exchange and explicitly targets active post-compromise settings; critically, within this model Signal’s user-mediated authentication cannot be proven secure, motivating a concrete fix intended to enable detection of an active MitM adversary \cite{dowling2021}. This line of work connects back to the broader secure-messaging framing that “secure” must incorporate properties beyond ciphertext confidentiality, and it exposes a tension between minimizing reliance on trusted third parties and achieving robust authentication when users are part of the trust loop \cite{oppliger2020,dowling2021}.

A third strand responds to the mismatch between two-party ratcheting designs and real-world group communication. Oppliger emphasizes that extending Signal-style guarantees to large groups is non-trivial and motivates MLS as a protocol built for thousands of participants, reflecting an architectural bifurcation: Signal’s design center is pairwise messaging, whereas MLS targets group scalability as a first-class requirement \cite{oppliger2024}. In parallel, Ayed and Dawood propose a protocol-design approach that retains Double Ratchet-inspired properties while engineering for group churn and asynchronous operation, notably by partitioning participants into capped subgroups treated as single entities for intra-group security management \cite{ayed2025}. Taken together, these contributions illustrate two distinct research strategies for group secure-messaging after 2020: standardization-driven redesign (MLS) versus adaptation of ratcheting principles via new group key agreement structures, each foregrounding different trade-offs between conceptual continuity with Signal and operational scalability \cite{oppliger2024,ayed2025}.

A fourth strand anticipates cryptographic migration pressures under quantum threat models, where secure-messaging must preserve Signal’s desirable operational properties while swapping out vulnerable primitives \cite{sharma2023,boras2025}. Sharma et al.\ argue that Signal’s reliance on ECDH (Curve25519) for key exchange becomes a critical weakness against quantum adversaries, motivating the replacement of classical key exchange with quantum-resistant schemes while retaining E2EE features such as forward secrecy and asynchronous communication \cite{sharma2023}. Boras’ literature review synthesizes the broader effort to develop and evaluate quantum-resistant primitives suitable for secure instant messaging, explicitly positioning the Signal protocol as a primary integration target \cite{boras2025}. Viewed alongside the group-messaging trajectory, this suggests that post-2020 secure-messaging work is increasingly “upgrade-oriented”: rather than proposing wholly new applications, it concentrates on how to evolve deployed protocols (Signal-like or MLS-like) without sacrificing their hard-won security and usability characteristics \cite{oppliger2024,boras2025}.

Finally, secure-messaging research has begun to more explicitly confront moderation, abuse handling, and sensitive-content risks as protocol-adjacent requirements that E2EE alone does not resolve \cite{jiang2024,samper2025}. Jiang et al.\ propose sMAC, a framework enabling sanitization and moderation over encrypted messages with formal guarantees including data privacy and sender anonymity, directly addressing the challenge that encrypted payloads impede abuse detection and accountability \cite{jiang2024}. At a higher level, Samper and Ferreira’s systematization of “safer sexting” technologies—spanning app features and academic proposals—reinforces that private chats in the real world demand a bundle of security, privacy, and accountability mechanisms, not a single cryptographic guarantee \cite{samper2025}. In synthesis, these governance-focused efforts complement authenticity and PQ migration work by broadening the secure-messaging agenda: modern secure-messaging is increasingly defined by how E2EE protocols interoperate with authentication UX, group scalability mechanisms, and privacy-preserving accountability or moderation layers \cite{dowling2021,oppliger2024,jiang2024,samper2025}.

\subsection{security-analysis}
Security analysis of Signal-protocol-based instant messaging since 2020 spans both implementation-layer compromise and protocol-/system-level weaknesses that persist even under end-to-end encryption. On the implementation side, Jain et al.\ demonstrate how threat modeling can translate into actionable exploit development against a widely deployed Signal-protocol messenger, culminating in an attack that subverted application security to access sensitive local user content (e.g., documents and media) despite the presence of E2EE \cite{jain2021}. This highlights a recurring security-analysis finding across modern messengers: confidentiality guarantees provided by the cryptographic channel do not, by themselves, prevent compromise via application ecosystem vulnerabilities and client-side attack surfaces \cite{jain2021}.

At the system level, Zhao et al.\ show an analogous disconnect in the 5G/RCS ecosystem: even when E2EE is deployed, weak cellular identity binding enables service hijacking that in turn permits impersonation and eavesdropping \cite{zhao2022}. Juxtaposed with Jain et al., these results emphasize that security analyses increasingly target the binding between identity, device/service provisioning, and cryptographic sessions, rather than treating E2EE as an end-state \cite{jain2021,zhao2022}. Both lines of work motivate defenses that extend beyond the core protocol to include stronger authentication/binding and hardening of surrounding infrastructure and clients \cite{jain2021,zhao2022}.

A complementary thread examines metadata privacy as a first-class security property. Brigham and Hopper analyze Signal’s Sealed Sender and Private Groups extensions and show that, for group conversations, communicating entities can be linked through recipient metadata alone, defeating the intended protection against a potentially compromised server \cite{brigham2025}. Their evaluation also stresses the limits of incremental defenses: a previously proposed mitigation for pairwise Sealed Sender traffic analysis does not carry over to the group setting, and the authors instead discuss a user-deployable (“server-agnostic”) mitigation with explicit overhead--efficacy tradeoffs \cite{brigham2025}. In contrast to compromise-oriented analyses that seek direct data access \cite{jain2021} or account/session takeover \cite{zhao2022}, this work underscores that security analysis in Signal-based messengers must also account for inferential attacks where message content remains encrypted but communication patterns leak \cite{brigham2025}.

Finally, several papers frame security analysis through protocol construction choices and performance/security tradeoffs. Cao and Zhang provide an algorithmic decomposition of Signal (including a mathematical walkthrough of X3DH) intended to support reasoning about how component cryptography composes into end-to-end transmission security, while also foregrounding performance considerations \cite{cao2022}. Angom et al.\ pursue a hardening direction for initial key establishment by proposing MLXDH, which integrates NIST’s ML-KEM and ML-DSA into Signal’s setup to add post-quantum hardness via PQ signatures, while acknowledging remaining discrete-log-based elements for mutual authentication in their design \cite{angom2025}. Although positioned as an enhancement rather than an attack study, this reflects a security-analysis emphasis on identifying residual assumptions (e.g., authentication dependencies) and incrementally raising the bar against future adversaries \cite{angom2025}. In parallel, Thapa et al.\ analyze encrypted messaging apps (including Signal) in terms of device resource consumption (CPU/GPU/memory) during encrypted transmission and propose CPU offloading to reduce overhead while maintaining acceptable confidentiality, illustrating how “security” in practice is analyzed alongside feasibility constraints on mobile hardware \cite{thapa2025}. Taken together, these works suggest that post-2020 security analysis around Signal and instant messaging is increasingly multi-layered: attackers and defenders alike must consider client compromise \cite{jain2021}, identity/service binding failures \cite{zhao2022}, metadata inference \cite{brigham2025}, and the evolving cryptographic and performance constraints shaping protocol and system design \cite{cao2022,angom2025,thapa2025}.

\subsection{security-audit}
Work on security auditing for post-2020 instant messaging interoperability has begun to focus on threats that arise \emph{outside} the end-to-end encryption layer—specifically, on detecting equivocation and order-integrity failures introduced by required intermediaries. In the IETF MIMI architecture for cross-platform group messaging, all MLS ciphertexts are routed through a central hub that timestamps and rebroadcasts messages; this design preserves confidentiality but creates a single point where a compromised hub can silently drop, delay, or reorder messages without breaking encryption guarantees \cite{sarvaiya2025}. 

Sarvaiya and Burger operationalize auditing as a lightweight, client-driven accountability layer that targets this hub-induced gap between cryptographic secrecy and verifiable delivery/order \cite{sarvaiya2025}. Their approach treats the hub’s timestamped broadcast stream as an auditable log: each client maintains an ordered list of received ciphertexts plus hub timestamps, then periodically emits Merkle proofs derived from its local list by embedding the proof within an encrypted application message \cite{sarvaiya2025}. This design choice links auditability to normal messaging traffic and makes the auditing cadence unpredictable to the hub, constraining a malicious hub’s ability to selectively suppress evidence-bearing messages \cite{sarvaiya2025}. The protocol-level consequence is a distributed cross-checking mechanism: once any proof is observed, other clients can validate it against their own histories and respond with their own proofs, so inconsistencies (from dropping or reordering) become visible to the room and accumulate into a non-repudiable record of hub misbehavior \cite{sarvaiya2025}. 

Within this theme, the contribution is less about strengthening the Signal/MLS cryptographic core than about furnishing an audit primitive suited to interoperability deployments where a shared relay is unavoidable \cite{sarvaiya2025}. The evaluation further positions audit as a practical add-on rather than a heavy consensus/log-replication system: in a Rust prototype built on OpenMLS and a 100-node emulated network, modest sampling (5\%) combined with a modeled hub attack probability (10\%) yielded high likelihood of early detection (95\% within the first 40 messages) at low per-client overhead (about 3\,kB memory and $<$1\,ms processing per proof) \cite{sarvaiya2025}. Collectively, this indicates an emerging audit direction for interoperable instant messaging since 2020: probabilistic, client-embedded verification layers that aim to expose intermediary misbehavior affecting delivery and ordering while keeping the end-to-end encryption substrate unchanged \cite{sarvaiya2025}.

\subsection{signal-app}
Since 2020, Signal-focused research in instant messaging has clustered around two tensions: (i) the extent to which end-to-end encryption meaningfully limits what can be learned from the system, and (ii) how the practical deployment of Signal on mobile devices creates residual attack and analysis surfaces. One strand demonstrates that, even when message contents are protected, Signal-related artifacts remain accessible in client-side storage under certain conditions, enabling recovery of communications for investigative purposes \cite{krishnapriya2021}. A second, complementary strand shows that encrypted network traffic itself can leak sensitive behavioral or social information, motivating both measurement infrastructures for Signal traffic and new defenses against metadata inference \cite{erdenebaatar2023,bozorgi2023,brigham2025}. A smaller line of work places Signal in comparative performance/security discussions, framing efficiency as an additional constraint shaping real-world protocol and app choices \cite{thapa2025}.

Across these papers, “what encryption protects” is scoped differently, producing a more complete view when combined. Device forensics work treats Signal’s local data as a primary evidentiary source, proposing acquisition via decryption of Signal backup files to recover messages and attachments on Android \cite{krishnapriya2021}. In contrast, traffic-analysis work assumes encryption holds and instead extracts information from observable packet/flow patterns, showing that sensitive user activities and relationships can be inferred merely by monitoring encrypted IM traffic across apps that include Signal \cite{bozorgi2023}. These perspectives are not contradictory but cumulative: they indicate that confidentiality of content does not preclude (a) recoverability from endpoint artifacts under certain access models \cite{krishnapriya2021} and (b) metadata/behavioral inference from encrypted traffic under passive observation models \cite{bozorgi2023}.

Methodologically, recent work also reflects a shift from ad hoc measurement to more systematic data generation for Signal traffic studies. The lack of publicly available encrypted IMA datasets motivates an automated traffic generation and labeling system that includes Signal and runs directly on Android, using emulated user behavior and on-device capture/filtering to produce analyzable traces \cite{erdenebaatar2023}. This kind of instrumentation underpins and accelerates the feasibility of the broader class of attacks demonstrated in traffic-analysis research (e.g., inferring participants in communications despite encryption) by enabling controlled, repeatable capture of app-specific encrypted traffic signatures \cite{erdenebaatar2023,bozorgi2023}. Together, these works suggest that progress on Signal traffic analysis depends not only on new inference techniques but also on the availability of reproducible, labeled datasets and test harnesses \cite{erdenebaatar2023,bozorgi2023}.

A central synthesis point is that privacy-preserving extensions in Signal are themselves active targets of analysis. While Signal’s Sealed Sender and Private Groups aim to conceal “who talks to whom” and group membership from potentially compromised servers, traffic-analysis research identifies ways to link group conversations through recipient metadata alone, effectively undermining these protections in the group setting \cite{brigham2025}. This result complements the broader finding that encrypted IM traffic can leak sensitive participation information without exploiting software vulnerabilities \cite{bozorgi2023}, collectively emphasizing that metadata protection is fragile even when cryptographic content protection is strong \cite{bozorgi2023,brigham2025}. Both works converge on countermeasure framing: one proposes general countermeasures and introduces a client-side system (IMProxy) that does not require provider support \cite{bozorgi2023}, while the other discusses a “server-agnostic” mitigation implementable by users that trades communication overhead for defensive strength \cite{brigham2025}. Read together, they highlight a pragmatic design space in which deployable, user-side defenses are pursued precisely because app/provider changes may be slow or unavailable \cite{bozorgi2023,brigham2025}.

Finally, beyond security and privacy leakage, Signal is examined as a concrete instantiation of end-to-end encryption with performance costs on mobile devices. Comparative measurement across Signal, WhatsApp, and Telegram evaluates CPU/GPU/memory footprints during encrypted messaging and motivates an optimization approach based on CPU offloading to reduce resource utilization while maintaining confidentiality expectations \cite{thapa2025}. This efficiency perspective connects back to both forensic and traffic-analysis lines: if resource constraints shape implementation choices and runtime behavior, they may indirectly affect observable traffic patterns and artifact management strategies, influencing both measurement and attack/defense evaluations in Signal deployments \cite{thapa2025,erdenebaatar2023,bozorgi2023}. Overall, post-2020 Signal-app research thus spans endpoint artifact extraction, systematic traffic dataset generation, and escalating scrutiny of metadata leakage (including within Signal’s own privacy extensions), alongside emerging attention to efficiency tradeoffs in encrypted messaging on constrained devices \cite{krishnapriya2021,erdenebaatar2023,bozorgi2023,brigham2025,thapa2025}.

\subsection{signal-based}
Signal-based research since 2020 in the provided set extends beyond messenger applications toward broader social networking architectures, but it diverges in how concretely Signal is operationalized. Basem et al. treat Signal as a protocol foundation to be systematically adapted for social network platforms, arguing that E2EE adoption has largely remained within messengers despite privacy risks on social platforms \cite{basem2023}. In contrast, Kanimozhi et al. embed ``end-to-end encryption similar to Signal'' as one component within a decentralized, AI-augmented social network design, positioning Signal-style E2EE as a baseline privacy mechanism alongside platform features such as local data ownership, ephemeral caching, and AI-driven interaction tooling \cite{kanimozhi2025}.

A key distinction is the depth of protocol redesign for social networking requirements. Stick is explicitly framed as a Signal-based protocol \emph{tailored} to the asynchronous, multi-device realities of social network platforms, introducing a new session concept, multiple pairwise sessions, and refreshing identity keys to support re-establishable encryption sessions while preserving forward secrecy and adding backward secrecy \cite{basem2023}. This focus on re-establishment and multi-device asynchrony foregrounds a protocol-level response to social-network usage patterns (e.g., many-to-many interactions and intermittent connectivity) rather than treating Signal-like encryption as a drop-in component \cite{basem2023}. By comparison, the decentralized social networking platform described by Kanimozhi et al. emphasizes system-level data placement (profiles and content stored locally, temporarily mirrored on decentralized nodes, erased post-session) and user experience features (e.g., ``Time Capsule'' messaging and in-chat agents), with Signal-like E2EE described at a higher level of abstraction \cite{kanimozhi2025}. Taken together, these works suggest two complementary trajectories: one that modifies Signal-derived mechanisms to meet social-network protocol demands, and another that treats Signal-style encryption as an enabling primitive within a broader privacy- and functionality-driven platform redesign \cite{basem2023,kanimozhi2025}.

The papers also contrast in their approach to security assurance and claimed security properties. Stick couples its Signal-based design with symbolic-model formal verification in Verifpal and presents a security analysis claiming a form of post-compromise security for many-to-many communications, alongside session re-establishment with authentication and confidentiality \cite{basem2023}. This verification-and-implementation pipeline (including an open-source API and performance/usability evaluation in a social-network app setting) reflects an emphasis on demonstrable correctness and deployability of Signal-derived protocol changes \cite{basem2023}. Kanimozhi et al., while likewise motivated by privacy and security, articulate trust and abuse-resistance largely through a ``quantum inspired'' reputation mechanism and AI-mediated deterrence in a decentralized architecture, with E2EE framed as Signal-like but not subjected to the same protocol-specific verification or property articulation within the abstract \cite{kanimozhi2025}. Across both, Signal remains a reference point for E2EE expectations, but only Stick operationalizes that reference into formally analyzed protocol mechanisms tailored to social-network constraints \cite{basem2023,kanimozhi2025}.

\subsection{signal-protocol}
Since 2020, research related to the Signal protocol in instant messaging has clustered into three intersecting strands: (i) formal and logic-based reasoning about Signal’s security goals, (ii) empirical and adversarial security analysis of Signal-based implementations, and (iii) forward-looking redesign efforts—especially post-quantum hardening—alongside consolidating explanatory resources that position Signal relative to emerging standards for group messaging \cite{oppliger2020,kamkuemah2021,jain2021,sharma2023,oppliger2024,boras2025,angom2025}.

A notable throughline is the use of formalization to make Signal’s guarantees explicit and checkable at different abstraction levels. Vladimirovich et al.\ model “signaling protocol” organization for online audio/video calls using non-deterministic event-driven automata, emphasizing how formal models can expose protocol factors and developer-facing problems in call setup flows \cite{vladimirovich2020}. Complementing this structural modeling, Kamkuemah applies epistemic logic to reason about authentication and secrecy in Signal in terms of “who knows what,” arguing that articulating such knowledge-based security properties before implementation supports security-by-design and concluding that the protocol is secure under the expressed properties \cite{kamkuemah2021}. In contrast to these formal and property-driven approaches, Cao and Zhang present a mathematically oriented decomposition of Signal’s X3DH calculation process as part of a broader comparison of end-to-end transmission schemes, using this breakdown to discuss how constituent algorithms compose into a protocol and to surface performance-related issues \cite{cao2022}. Together, these works reflect a shift from treating Signal as a black-box “secure messenger” primitive toward making its protocol mechanics and desired guarantees explicit in models and stepwise derivations, albeit with differing goals (developer clarity for calls vs.\ logic-level secrecy/authentication vs.\ algorithmic decomposition and performance discussion) \cite{vladimirovich2020,kamkuemah2021,cao2022}.

Running alongside formal reasoning is work that stresses the gap between protocol-level assurances and application/ecosystem realities. Jain et al.\ analyze a large-scale consumer messenger that uses the Signal protocol and, after threat modeling and executing attacks, report an attack that subverted the application’s security to access confidential user data (e.g., documents and media), reproducing the attack across systems and reporting findings to authorities \cite{jain2021}. This contrasts with logic-based conclusions of protocol security by highlighting that vulnerabilities can arise from the broader ecosystem or implementation decisions even when the underlying protocol is argued to satisfy key security properties \cite{kamkuemah2021,jain2021}. In this sense, the post-2020 literature reflects a two-level security narrative: protocol analyses aim to validate secrecy/authentication and related properties in an abstract model, while vulnerability analyses demonstrate that end-user confidentiality can still fail due to weaknesses outside (or adjacent to) the cryptographic core \cite{kamkuemah2021,jain2021}. Survey-style overviews also contribute to this framing by presenting Signal as a core enabler of end-to-end encrypted messaging while explicitly drawing attention to broader privacy and metadata issues (e.g., traffic analysis) that persist beyond message content confidentiality \cite{oppliger2020}.

A second major trajectory since 2023 is post-quantum adaptation of Signal’s key establishment and related primitives. Sharma et al.\ motivate replacing Signal’s ECDH Curve25519 key exchange due to quantum threats, positioning Signal’s current feature set (forward secrecy, future secrecy, authentication, asynchronous communication, deniability) as a baseline that post-quantum alternatives should preserve \cite{sharma2023}. Boras’ literature review similarly frames the post-quantum problem as a need to integrate quantum-resistant primitives into secure instant messaging protocols, with an explicit focus on Signal as a target for such transition efforts \cite{boras2025}. Moving from motivation and survey to specific construction, Angom et al.\ propose MLXDH as a proof-of-concept hardening of Signal’s initial key establishment (building on PQXDH) using NIST ML-KEM and ML-DSA; they argue that while mutual authentication remains discrete-logarithm based in their design, adding post-quantum signatures increases cryptographic hardness and they discuss routes toward fully post-quantum-secure mutual authentication \cite{angom2025}. Across these papers, the post-quantum thread converges on initial key establishment as the primary pressure point for quantum resilience, but differs in maturity: from problem framing and requirements anchored in Signal’s existing guarantees, to a literature-level mapping of candidate primitives, to a concrete hybridized design and implementation sketch that strengthens parts of the handshake while acknowledging remaining classical dependencies \cite{sharma2023,boras2025,angom2025}.

Finally, multiple overview resources since 2020 consolidate and contextualize Signal protocol knowledge, while linking it to the broader evolution of secure messaging protocols and scalability challenges. Oppliger’s 2020 book situates Signal among other E2EE technologies and messengers, explains WhatsApp’s use of Signal, and foregrounds privacy issues such as metadata-driven traffic analysis—implicitly underscoring that “Signal protocol adoption” does not exhaust the privacy problem \cite{oppliger2020}. Oppliger’s 2024 volume deepens the protocol exposition and explicitly couples Signal to the challenge of scaling E2EE from two-party messaging to large groups, using Messaging Layer Security (MLS) as the standardization-oriented response for large-scale group messaging and positioning the book as a bridge between Signal’s mechanisms and MLS’s group-focused design goals \cite{oppliger2024}. In parallel, shorter survey accounts (e.g., Vidya Sagar et al.) frame Signal’s encryption in high-level public/secret key terms as part of a broader discussion of E2EE messaging security \cite{p2023}. Read together, these resources suggest that post-2020 work has increasingly treated Signal not merely as an app-specific protocol but as a reference point for (a) systematic explanation and practitioner education, (b) comparative protocol analysis and performance discussion, and (c) evolution toward new standards for group settings and new threat models such as quantum adversaries \cite{oppliger2020,cao2022,oppliger2024,boras2025}.

Overall, research since 2020 shows Signal protocol scholarship expanding from descriptive and comparative accounts into more explicit formal reasoning and adversarial testing, while simultaneously preparing for protocol evolution under post-quantum constraints and large-group scalability pressures \cite{vladimirovich2020,kamkuemah2021,jain2021,oppliger2024,angom2025}.

\subsection{signal-protocol-components}
Work since 2020 engaging Signal-protocol components splits between (i) performance-motivated reconfigurations of the X3DH + Double Ratchet stack and (ii) formal scrutiny of how Signal’s surrounding authentication workflow composes with ratcheting. On the engineering side, an ORDEX-based end-to-end encrypted chat design keeps the familiar Signal-style structure—initial key establishment via X3DH and ongoing key evolution via the Double Ratchet—but swaps in an ORDEX-derived key exchange intended to reduce the overhead attributed to conventional ECDH-based exchanges \cite{m2024}. In this view, Signal components are treated as modular building blocks that can be paired with alternative key exchange mechanisms while retaining the same high-level message-encryption pipeline (with AES used for message confidentiality) \cite{m2024}.

In contrast, formal methods work argues that security of a ratcheted design cannot be assessed solely by its confidentiality properties and internal component choices; instead, it depends critically on the entity-authentication layer that is often user-mediated and out-of-band \cite{dowling2021}. Dowling and Hale propose a model for user-mediated entity authentication that is explicitly designed to compose with any ratcheted key exchange, thereby targeting systems that adopt Signal-like ratcheting components \cite{dowling2021}. Applying this model to Signal’s user-mediated authentication, they find it cannot be proven secure against active adversaries in a post-compromise setting and propose a simple adjustment that enables detection of active MitM behavior \cite{dowling2021}. Taken together, these directions highlight a tension in post-2020 research on Signal-protocol components: while component modularity encourages swapping primitives to optimize performance (e.g., replacing ECDH-oriented exchange with ORDEX while retaining X3DH and Double Ratchet), correctness and security must also account for how authentication rituals and user checks interact with ratcheted exchanges under active compromise \cite{m2024,dowling2021}.

\subsection{survey-or-overview}
Survey and overview work since 2020 positions the Signal protocol as the de facto reference point for modern instant messaging end-to-end encryption, while increasingly framing it within broader ecosystems and adjacent standards. Oppliger’s book-length treatment situates Signal alongside earlier secure-messaging approaches and explicitly connects it to real-world deployments (e.g., WhatsApp) and privacy concerns beyond message content, such as metadata-enabled traffic analysis \cite{oppliger2020}. This broader scoping contrasts with more protocol-centric overviews that treat Signal primarily as a composition of cryptographic building blocks: Cao and Zhang decompose Signal (alongside BitTorrent and a smart city framework) into constituent algorithms and include a mathematically detailed walk-through of X3DH, using this decomposition to foreground how protocol construction impacts security and performance considerations \cite{cao2022}. In parallel, application-oriented survey discourse often abstracts Signal into a general “secure chat” pattern—public/private key usage and hybrid encryption framing—presenting Signal as an integration of multiple cryptographic methods aimed at mitigating common data-attack threats in messaging \cite{p2023}. Together, these sources show a spectrum of survey granularity: from ecosystem/implementation and privacy framing \cite{oppliger2020}, to algorithmic decomposition and handshake explication \cite{cao2022}, to higher-level secure-chat narratives \cite{p2023}.

A key throughline in post-2020 overviews is the tension between Signal’s two-party origins and the practical demands of group messaging at scale. Oppliger’s later work explicitly pairs Signal with Messaging Layer Security (MLS), presenting MLS as the IETF-standardizing answer to scaling end-to-end encrypted group communication to very large groups, and thereby positioning Signal not only as a “gold standard” baseline but also as a foil that motivates MLS’s design goals \cite{oppliger2024}. This coupling reframes “Signal-related” survey work from describing a single protocol to mapping a design space in which Signal is one end of the continuum (two-party E2EE guarantees) and MLS targets the operational realities of large-group deployments \cite{oppliger2024}. In this sense, the survey literature since 2020 uses Signal both as a canonical protocol to explain and as a benchmark against which emergent group-messaging solutions are introduced \cite{oppliger2024}.

Another integrative trend is the expansion from protocol security to practical efficiency and deployment trade-offs in consumer messengers. Thapa et al.\ compare WhatsApp, Telegram, and Signal by measuring resource consumption (CPU, GPU, memory) during encrypted message transmission, treating Signal-style E2EE as desirable but potentially costly on constrained mobile hardware \cite{thapa2025}. While Oppliger’s 2020 overview emphasizes privacy issues such as metadata and traffic analysis that persist even with E2EE \cite{oppliger2020}, Thapa et al.\ highlight a different “beyond-confidentiality” consideration—runtime overhead—and propose CPU offloading as an optimization intended to preserve confidentiality while reducing device burden \cite{thapa2025}. Read together, these overviews imply that the “Signal protocol in practice” is evaluated not only in cryptographic terms but also through systems criteria (resource usage) that shape user experience and feasibility of secure defaults \cite{thapa2025}.

Finally, survey-oriented work that is not protocol-specific still leverages instant messaging as a primary threat surface and design context, indirectly reinforcing the importance and limits of E2EE baselines like Signal. Samper and Ferreira’s SoK on safer sexting reviews a wide set of technology proposals and real-world app features across messaging and social platforms, emphasizing that safety and privacy in private chats extends beyond any single technique and requires threat-model-driven feature taxonomies \cite{samper2025}. This broader synthesis complements Signal-centered overviews by underscoring that even if Signal-like E2EE secures message content, users face additional risks and objectives (privacy, accountability, harm reduction) that are addressed through a diverse ecosystem of features rather than an “ultimate solution” \cite{samper2025}. Collectively, survey and overview research since 2020 thus treats Signal as foundational for confidential messaging \cite{oppliger2020,oppliger2024}, while increasingly embedding it within larger considerations: protocol composition and performance \cite{cao2022}, usability- and application-level framings of secure chat \cite{p2023}, mobile efficiency constraints \cite{thapa2025}, and broader private-chat safety taxonomies \cite{samper2025}.

\subsection{system-design}
System-design work since 2020 spans both communication-layer efficiency and end-to-end security architecture, with different papers repositioning where computation and trust should reside in instant messaging systems. At the networking layer, eDirect reframes “always-online” presence as a systems problem driven by periodic heartbeats: by shifting uplink responsibility to opportunistically selected device-to-device (D2D) relays that aggregate many users’ heartbeat messages before contacting the base station, it targets cellular signaling storms and handset energy drain simultaneously \cite{yi2020}. This design treats IM reliability as dependent on infrastructure-facing behavior (channel establishment/release overhead) and introduces scheduling based on heartbeat deadlines and frequencies to preserve timeliness while reducing signaling and energy costs \cite{yi2020}.

By contrast, the 2025 designs concentrate on end-to-end encrypted messaging’s trust boundaries, particularly where encryption/decryption and authentication should occur. ARSecure explicitly reacts to the threat model of client-side scanning (CSS) by moving cryptographic operations off the potentially compromised client device and onto augmented reality glasses positioned as a secure external device \cite{alsop2025}. In doing so, it systemically re-partitions the endpoint: rather than attempting to “harden” the phone app against scanning or tampering, ARSecure’s design goal is tamper-resistant message handling that renders CSS-style on-device inspection less effective by avoiding processing plaintext on the conventional endpoint \cite{alsop2025}. This hardware-mediated approach complements, rather than replaces, conventional E2EE semantics by changing the locus of trusted computation and aiming for “tamper-proof, private communication” in sensitive contexts \cite{alsop2025}.

EnConvo tackles a different system-design pressure: long-term cryptographic resilience and operational assurance under quantum-era assumptions. Its architecture couples a post-quantum key exchange (CRYSTALS-Kyber-1024) with X25519 for backward compatibility, then relies on hardware-based AES-256 in OCB mode for data encryption and uses hardware enclaves (Intel SGX, ARM TrustZone) for key protection \cite{r2025}. Notably, EnConvo’s system design links protocol selection to verifiability and performance claims by adopting the Noise Protocol Framework and using ProVerif to check correct execution, while also reporting measured latencies for the Kyber exchange and per-message encryption \cite{r2025}. Compared to ARSecure’s external-device split, EnConvo keeps computation on general-purpose devices but attempts to reduce the trusted software base through enclaves and zero-knowledge authentication (zk-SNARKs) for passwordless login without revealing credentials \cite{r2025}. Across both, the “endpoint” is no longer a monolithic app: ARSecure externalizes it into an AR device to evade client-side compromise \cite{alsop2025}, while EnConvo internalizes protections via enclaves and ZK authentication to constrain credential and key exposure \cite{r2025}.

A third design direction broadens the system boundary beyond a single messaging app into a decentralized social networking platform that borrows Signal-like end-to-end encryption while rethinking data custody and availability. Kanimozhi et al.\ propose local-first storage of profiles and content on user devices paired with ephemeral caching servers and decentralized nodes that only mirror data during active interactions and delete it afterward \cite{kanimozhi2025}. This architecture implicitly trades continuous server-side persistence for session-scoped replication, aligning availability with interaction windows rather than platform-wide retention \cite{kanimozhi2025}. While EnConvo and ARSecure focus on cryptographic and endpoint trust, this platform-level design incorporates additional system components (AI personalization, group agents, “Time Capsule” unlocking, and a trust-score mechanism) that sit adjacent to messaging and therefore shape how encrypted communication is embedded into broader social features \cite{kanimozhi2025}. Importantly, its explicit aspiration to be “similar to Signal” in E2EE indicates that Signal-protocol ideas are serving as a design reference point even when the surrounding system is decentralized and augmented with AI-mediated functions \cite{kanimozhi2025}.

Taken together, these system-design efforts highlight a diversification in post-2020 instant messaging research: efficiency-focused redesigns of background connectivity to reduce infrastructure and battery costs \cite{yi2020}, endpoint re-partitioning to mitigate client-side inspection and tampering risks \cite{alsop2025}, cryptographic agility with post-quantum and formally checked protocol stacks paired with hardware key protection \cite{r2025}, and platform architectures that combine Signal-like E2EE with decentralized, ephemeral data replication and AI-driven features \cite{kanimozhi2025}. Rather than converging on a single “best” system architecture, the papers collectively suggest that IM system design is being pulled by multiple constraints—cellular signaling/energy, endpoint integrity under CSS, post-quantum preparedness, and data sovereignty—each motivating different placements of state, computation, and trust \cite{yi2020,alsop2025,r2025,kanimozhi2025}.

\subsection{traffic-analysis}

Across post-2020 work, traffic analysis emerges as the primary lens through which researchers interrogate what encrypted instant messaging (including Signal and the Signal protocol ecosystem) continues to leak beyond message contents. The literature collectively positions traffic analysis as a metadata problem: even when end-to-end encryption is sound, observable communication patterns (e.g., timing, volume, and recipient-side information) can still enable inference about users and their interactions \cite{oppliger2020,bozorgi2023,brigham2025}. In this framing, encryption shifts the adversarial focus from cryptographic breakage to side-channel inference on network traces and protocol-visible metadata \cite{oppliger2020,bozorgi2023}.

A key line of research demonstrates that encrypted traffic can support high-stakes inferences about participation in sensitive communications without exploiting application vulnerabilities \cite{bozorgi2023}. Bozorgi et al.\ show that monitoring encrypted IM traffic alone can enable an adversary to identify participants in target communications (including forum-like settings) with high accuracy, emphasizing that these inferences are feasible in realistic experimental conditions and therefore constitute a practical surveillance risk \cite{bozorgi2023}. This attack-centric work complements Oppliger's broader discussion that metadata can be used for traffic analysis against E2EE messengers (including those built around the Signal protocol), underscoring that privacy risks persist even when content confidentiality holds \cite{oppliger2020}.

More recent Signal-specific analysis refines the threat model to include mechanisms explicitly designed to hide metadata from the server. Brigham and Hopper analyze Signal group conversations under Sealed Sender and Private Groups and argue that groups of communicating entities can still be linked using recipient metadata alone, defeating the intended protections in those extensions \cite{brigham2025}. Conceptually, this result aligns with the broader theme that encryption and even metadata-hiding features may leave residual, structurally exploitable signals; however, it shifts attention from general encrypted-traffic patterns to how protocol- and feature-level metadata exposure can re-enable linkage at the group level \cite{brigham2025}. Their evaluation also stresses that defenses proposed for related pairwise settings do not necessarily generalize to the group case, suggesting that countermeasure design must be sensitive to communication topology (pair vs.\ group) and the specific metadata surfaces available \cite{brigham2025}.

Alongside attack demonstrations, the literature highlights the practical bottleneck of obtaining labeled encrypted-traffic datasets for IMAs, which directly affects the reproducibility and breadth of traffic-analysis research \cite{erdenebaatar2023}. Erdenebaatar et al.\ respond by designing an encrypted traffic generation system that automates user-behavior emulation, on-device capture, filtering, and labeling across multiple IMAs including Signal, thereby supporting data-driven encrypted traffic analysis where public datasets are scarce \cite{erdenebaatar2023}. This dataset-generation emphasis complements attack papers by supplying the experimental substrate needed to train and validate flow-based or statistical inference approaches at scale \cite{erdenebaatar2023,bozorgi2023}. Taken together, these works suggest a research pipeline in which the feasibility of sensitive inference from encrypted traces \cite{bozorgi2023,brigham2025} motivates both systematic data generation \cite{erdenebaatar2023} and renewed attention to metadata-centric privacy mitigations in E2EE messaging \cite{oppliger2020}.

Mitigation discussions further reveal a divide between solutions requiring provider integration versus user-deployable approaches. Bozorgi et al.\ report that standard countermeasures can reduce attack effectiveness and propose IMProxy as a client-side countermeasure that does not require support from IM providers, directly engaging the deployment friction that often slows protective changes in messaging ecosystems \cite{bozorgi2023}. Brigham and Hopper similarly explore a ``server-agnostic'' mitigation implementable by users, but explicitly characterize a tradeoff between added communication overhead and defense efficacy in the group setting \cite{brigham2025}. In synthesis, recent work converges on the view that traffic-analysis resistance is not a binary property provided by encryption alone: it is an evolving systems problem where (i) observable traffic or feature-level metadata can still permit linkage and participation inference \cite{bozorgi2023,brigham2025}, (ii) rigorous evaluation depends on high-quality labeled captures of encrypted app traffic \cite{erdenebaatar2023}, and (iii) pragmatic mitigations increasingly aim to be user-deployable while managing overhead and usability constraints \cite{bozorgi2023,brigham2025}.

\subsection{vulnerability-analysis}
Since 2020, vulnerability-analysis work at the intersection of instant messaging and modern secure messaging deployments has emphasized that end-to-end encryption (including Signal-derived designs) does not, by itself, close off compromise paths that arise from application ecosystems and service-layer identity binding \cite{jain2021,zhao2022}. Across two different settings—an over-the-top Signal-protocol-based consumer messenger and carrier-mediated 5G RCS—both studies adopt threat-driven methodologies to identify weaknesses that persist even when message payloads are cryptographically protected \cite{jain2021,zhao2022}.

A key shared insight is the centrality of identity and endpoint assumptions as practical attack surfaces. Jain et al.\ ground their analysis in threat modeling of a Signal-protocol-based application and then execute attacks that exploit identified weaknesses, including an attack that “successfully subverted the application’s security” to access confidential user data such as documents and media \cite{jain2021}. In contrast, Zhao et al.\ focus on the messaging system architecture of 5G RCS and empirically validate attacks across major mobile carriers; they attribute feasibility of impersonation and eavesdropping to weak cellular ID binding, noting that such attacks remain possible “even with end-to-end encryption in place” \cite{zhao2022}. Taken together, these works converge on the idea that identity binding and endpoint/platform integration can dominate risk, so that encryption at the transport/content layer becomes necessary but insufficient for user security outcomes \cite{jain2021,zhao2022}.

Methodologically, the papers illustrate complementary approaches to vulnerability analysis that span application-level and infrastructure-level scopes. Jain et al.\ combine a schema-level understanding “based on publicly available knowledge” with threat modeling frameworks, then reproduce attacks on multiple systems and report findings to authorities, reflecting a workflow oriented around actionable exploitation evidence within a single widely deployed app ecosystem \cite{jain2021}. Zhao et al.\ provide an ecosystem-wide security study for a standardized messaging service, validating attacks in-the-wild across carriers under ethical constraints and pairing the attack study with proposed immediate and long-term defenses, reflecting a workflow oriented toward systemic remediation across stakeholders \cite{zhao2022}. The juxtaposition suggests that post-2020 vulnerability analysis in messaging increasingly spans (i) app/OS/file-system and client-ecosystem exposure in Signal-protocol-based messengers and (ii) service identity, provisioning, and cross-device assumptions in carrier-grade messaging, with both strands using empirical attack validation to move beyond purely theoretical cryptographic analysis \cite{jain2021,zhao2022}.

\subsection{x3dh}
Since 2020, work referencing X3DH in the instant-messaging context in this set clusters around two distinct directions: performance-motivated re-engineering of the initial handshake and security-hardening of the same phase toward post-quantum resilience. Both papers treat X3DH as the anchoring mechanism for initial shared-secret establishment before transitioning to ongoing message protection (e.g., via Double Ratchet), but they diverge sharply in what they change and what properties they prioritize \cite{m2024,angom2025}.

One line of work integrates X3DH into alternative, ostensibly faster key-exchange pipelines. Vishwas et al.\ propose combining an OR Diffie-Hellman Key Exchange (ORDEX) algorithm with X3DH and Double Ratchet, with AES used for message encryption, explicitly framing X3DH as part of a multi-layer design intended for organizational chat applications \cite{m2024}. Their core contribution is not a modification of X3DH’s authentication structure per se but an attempt to replace or augment conventional Diffie–Hellman-style operations with ORDEX to reduce overhead, arguing the resulting solution is “thousand times faster” than approaches using ECDH key exchange \cite{m2024}. In this view, X3DH functions as a composable component within a performance-optimized stack, where speed is the driving metric and X3DH is paired with Double Ratchet to sustain secure messaging after session setup \cite{m2024}.

A contrasting direction focuses on hardening X3DH-style initial establishment against future cryptanalytic threats, particularly quantum adversaries, while preserving Signal-like semantics. Angom et al.\ analyze Signal’s initial key establishment (discussed via PQXDH) and present MLXDH, a proof-of-concept construction that incorporates NIST’s ML-KEM and ML-DSA into the initial establishment process \cite{angom2025}. Rather than claiming performance gains, they motivate their design by increasing cryptographic hardness: although mutual authentication in their design remains discrete-logarithm based, they argue adding post-quantum signatures enhances security and they outline generic pathways toward a fully post-quantum-secure key establishment including mutual authentication \cite{angom2025}. Relative to the performance-driven ORDEX+X3DH integration \cite{m2024}, MLXDH treats the X3DH-family exchange as a security boundary to be strengthened with additional primitives, explicitly targeting post-quantum security properties while acknowledging remaining dependence on discrete-log assumptions for authentication \cite{angom2025}.

Taken together, these papers illustrate that recent X3DH-related work in messaging systems can be read as a tension between (i) engineering-oriented recomposition of the initial exchange to reduce perceived computational cost while keeping the familiar X3DH $\rightarrow$ Double Ratchet flow \cite{m2024}, and (ii) protocol-design extensions that preserve the initial exchange’s role but “wrap” it with post-quantum mechanisms to hedge against quantum threats, even if complete post-quantum mutual authentication remains an open design target in the proposed construction \cite{angom2025}.

\section{Conclusion}

Since 2020, research related to instant messaging and the Signal protocol has largely treated Signal as a de facto baseline for end-to-end encrypted (E2EE) communication and then explored where that baseline breaks down or must be extended. Rather than proposing wholesale replacements for Signal, much of the work either (i) stress-tests Signal-style security under more realistic system and user assumptions, or (ii) adapts Signal-derived components (notably X3DH and Double Ratchet) to emerging requirements such as post-quantum migration, large-group communication, interoperability, and multi-device use \cite{oppliger2020,oppliger2024,basem2023,angom2025,bobrysheva2021,ayed2025}. In parallel, a substantial thread broadens “secure messaging” beyond protocol cryptography to end-to-end exposure, showing that privacy and security failures often arise from endpoints, authentication practices, metadata leakage, and platform artifacts rather than from weaknesses in the core content-encryption design \cite{jain2021,krishnapriya2021,bozorgi2023,brigham2025}.

Across themes, the literature converges on a recurring contrast: content confidentiality is comparatively well addressed by Signal-style E2EE, while identity/authentication, metadata privacy, and endpoint integrity remain persistent weak points. Formal and logic-based analyses can validate secrecy/authentication properties under specific models \cite{kamkuemah2021}, yet more user-centered and post-compromise threat models show that Signal’s user-mediated key verification resists proof against active MitM adversaries and motivates protocol-level detection fixes \cite{dowling2021}. This authentication pressure connects directly to governance and trust distribution: decentralized-identity proposals argue that even perfect E2EE leaves the public-key authentication problem unresolved when providers remain de facto trust anchors, motivating DID/VC-based approaches that externalize authentication away from IM services and align with interoperability/regulatory drivers \cite{morales2024}. In other words, where Signal traditionally assumes a provider-mediated identity directory plus optional user verification, recent work increasingly treats authentication as the central socio-technical bottleneck rather than an auxiliary feature \cite{dowling2021,morales2024}.

A second cross-cutting relationship links “what Signal encrypts” with “what Signal leaks.” Traffic-analysis and metadata-focused studies demonstrate that encrypted traces can still reveal participation patterns, and that even Signal-specific protections (e.g., Sealed Sender/Private Groups) may allow group linkability via recipient metadata \cite{bozorgi2023,brigham2025}. Complementing these empirical findings, protocol-design and formal-modeling efforts attempt to add new privacy properties (e.g., deniable, metadata-private traffic) while preserving deployability by layering on top of unmodified Signal, as in DenIM \cite{nelson2024}. This pairing highlights an emerging design pattern: empirical leakage studies motivate new formalized privacy goals, and subsequent protocol variants attempt to achieve them with minimal disruption to the deployed ecosystem \cite{bozorgi2023,nelson2024,brigham2025}.

A third thematic bridge concerns scaling and integration pressures. Group messaging and interoperability research shows that moving from two-party Signal roots to large groups and cross-provider delivery introduces integrity and ordering risks even when confidentiality is preserved (e.g., hub drop/reorder in MIMI/MLS), prompting auditing mechanisms such as Merkle-proof sampling to make relay misbehavior detectable \cite{sarvaiya2025,oppliger2024}. Meanwhile, other group-oriented work extends ratcheting concepts or introduces subgroup partitioning and post-quantum group key establishment, reflecting a tension between adopting MLS as a standardized group-native solution and extending Signal-like ratcheting to satisfy bespoke requirements \cite{bobrysheva2021,ayed2025,oppliger2024}. Multi-device work such as Stick further underscores that “Signal-like” security properties must be re-argued when session continuity, device churn, and asynchronous many-to-many communication become first-class constraints, motivating explicit re-establishment mechanisms and symbolic verification as part of protocol engineering \cite{basem2023}.

Post-quantum (PQ) migration is another major driver that binds protocol-component research, surveys, and system design. Multiple efforts focus on hardening Signal’s initial key establishment (the X3DH stage) with NIST PQ primitives or hybrid approaches (e.g., MLXDH), while acknowledging that authentication may still rely on classical assumptions unless signatures and identity binding are also migrated \cite{sharma2023,angom2025,boras2025}. System-level designs extend this “stack” perspective by combining PQ key exchange with enclaves, zero-knowledge authentication, and formal checking, signaling a move from protocol-only upgrades to integrated end-to-end architectures \cite{r2025}. This PQ thread complements the authentication and metadata threads: strengthening confidentiality against future adversaries does not automatically address who is being authenticated, what metadata is exposed, or what happens under endpoint compromise \cite{morales2024,bozorgi2023,jain2021}.

Notably, several themes emphasize that Signal’s security is frequently circumvented at the endpoints and platforms where it is used. Forensics and data-extraction work focuses on decrypting Signal backups on Android to recover messages and attachments, demonstrating that evidentiary recovery and adversarial acquisition can proceed via at-rest artifacts without attacking in-transit cryptography \cite{krishnapriya2021}. Client-side scanning proposals and critiques similarly reframe the threat: if the client can be compelled or compromised to scan plaintext, E2EE guarantees become moot, motivating architectural repartitioning such as offloading encryption/decryption to dedicated external devices \cite{alsop2025}. Vulnerability analyses in Signal-based or E2EE-enabled ecosystems reinforce this point by demonstrating compromise routes that expose sensitive local data or enable persistent impersonation when identity binding is weak, even if message encryption remains intact \cite{jain2021,zhao2022}. Together, these lines indicate that post-2020 research has shifted from “is the protocol cryptographically sound?” toward “does the full system deliver end-to-end security under realistic compromise and governance conditions?” \cite{jain2021,alsop2025,krishnapriya2021}.

Despite this breadth, important gaps remain. First, there is a recurring disconnect between formal assurances and deployment realities: proofs and symbolic verification often assume idealized user behavior, device integrity, or authentication channels, while empirical and user-mediated models highlight precisely those weak assumptions as failure points \cite{kamkuemah2021,dowling2021,basem2023}. More work is needed to unify formal models with usable security constraints—especially for key verification, post-compromise MitM detection, and cross-device identity management—so that provable properties align with how Signal-like systems are actually used \cite{dowling2021,basem2023}. Second, metadata privacy remains only partially addressed: traffic-analysis attacks continue to scale, mitigations can impose overhead on users or infrastructures, and protocol variants that improve metadata privacy (e.g., DenIM) require broader evaluation of deployability, interoperability interactions, and adversaries that combine network observation with endpoint compromise \cite{bozorgi2023,brigham2025,nelson2024}. Third, PQ upgrades are still incomplete at the system level: many proposals strengthen key agreement but leave open questions about migrating identity binding and authentication, managing hybrid transitions, and validating performance and usability impacts in real mobile deployments \cite{angom2025,sharma2023,boras2025}. Fourth, the literature set reflects a platform skew: Android forensics is well represented, but comparable post-2020 work on iOS artifacts, desktop clients, and cross-platform backup behaviors is not visible here, limiting conclusions about endpoint exposure beyond Android backups \cite{krishnapriya2021}. Finally, interoperability and moderation/accountability directions point to unresolved tensions: auditing layers can improve integrity in hub-mediated systems \cite{sarvaiya2025}, and encrypted moderation frameworks aim to add governance without forfeiting privacy \cite{jiang2024}, but open questions remain about how these features compose with Signal-like privacy goals, decentralized identity, and user expectations in real-world deployments \cite{jiang2024,morales2024,sarvaiya2025}.

Overall, the post-2020 research landscape suggests that “research related to Signal” is increasingly research about the boundaries of Signal’s guarantees: extending Signal-derived mechanisms to new settings (groups, multi-device, interoperability), hardening them for future cryptographic threats (PQC), and confronting what E2EE does not solve (authentication trust, metadata leakage, endpoint compromise, and governance). The central implication is that the next advances are unlikely to come from changing message encryption alone; instead, progress depends on integrating authentication reform, metadata protections, integrity/auditability in interoperable deployments, and endpoint-resilient architectures into a cohesive, evaluable secure-messaging stack \cite{dowling2021,morales2024,brigham2025,sarvaiya2025,alsop2025,r2025}.

\printbibliography
\end{document}
